% Template for PLoS
% Version 3.8 Apr 2026
%
% % % % % % % % % % % % % % % % % % % % % %
%
% -- IMPORTANT NOTE
%
% This template contains comments intended 
% to minimize problems and delays during our production 
% process. Please follow the template instructions
% whenever possible.
%
% % % % % % % % % % % % % % % % % % % % % % % 
%
% Once your paper is accepted for publication, 
% PLEASE REMOVE ALL TRACKED CHANGES in this file 
% and leave only the final text of your manuscript. 
% PLOS recommends the use of latexdiff to track changes during review, as this will help to maintain a clean tex file.
% Visit https://www.ctan.org/pkg/latexdiff?lang=en for info.
%
%
% IMPORTANT
% Below are a few tips to help format manuscript according to our specifications. For more tips to help reduce the possibility of formatting errors during conversion, please see our LaTeX guidelines at http://journals.plos.org/plosone/s/latex
%
% There are no restrictions on package use within the LaTeX files except that no packages listed in the template may be deleted.
%
% Please do not include colors or graphics in the text.
%
% The manuscript LaTeX source should be contained within a single file (do not use \input, \externaldocument, or similar commands).
%
% % % % % % % % % % % % % % % % % % % % % % %
%
% -- FIGURES AND TABLES
%
% Please include tables/figure captions directly after the paragraph where they are first cited in the text.
%
% DO NOT INCLUDE GRAPHICS IN YOUR MANUSCRIPT
% - Figures should be uploaded separately from your manuscript file. 
% - Figures generated using LaTeX should be extracted and removed from the PDF before submission. 
% - Figures containing multiple panels/subfigures must be combined into one image file before submission.
% For figure citations, please use "Fig" instead of "Figure".
% See http://journals.plos.org/plosone/s/figures for PLOS figure guidelines.
%
% Tables should be cell-based and may not contain:
% - spacing/line breaks within cells to alter layout or alignment
% - do not nest tabular environments (no tabular environments within tabular environments)
% - no graphics or colored text (cell background color/shading OK)
% See http://journals.plos.org/plosone/s/tables for table guidelines.
%
% For tables that exceed the width of the text column, use the adjustwidth environment as illustrated in the example table in text below.
%
% % % % % % % % % % % % % % % % % % % % % % % %
%
% -- EQUATIONS, MATH SYMBOLS, OTHER SPECIAL FORMATTING
%
% - For bold symbols, use the \boldsymbol command and not \mathbf. 
% - For inline equations, please be sure to include all portions of an equation in the math environment.  For example, x$^2$ is incorrect; this should be formatted as $x^2$ (or $\mathrm{x}^2$ if the romanized font is desired).
% - Do not include text that is not math in the math environment. For example, CO2 should be written as CO\textsubscript{2} instead of CO$_2$.
% - Avoid capturing simple text as equations, for example $<$.
% - For inline equations, please do not include punctuation (commas, etc) within the math environment unless this is part of the equation.
% - When adding superscript or subscripts outside of brackets/braces, please group using {}.  For example, change "[U(D,E,\gamma)]^2" to "{[U(D,E,\gamma)]}^2". 
% - Do not use \cal for caligraphic font.  Instead, use \mathcal{}.
% - Avoid splitting for fragmenting a single equation into multiple objects/environments, for example $f(x)$ = $g(x)$ + $h(x)$.
% - Use math environment for all equations, do not mimic using text.
% - Avoid using single letters in math mode where possible.
% - Always use matching parentheses and delimiters, avoid mixing math and text brackets.
% - Insert spaces around inline equations to ensure proper display after conversion.
% - Do not use forced numbers for display equation labeling or suppress numbering with \nonumber. 
%
% % % % % % % % % % % % % % % % % % % % % % % % 
%
% Please contact customercare@plos.org with any questions.
%
% % % % % % % % % % % % % % % % % % % % % % % %

\documentclass[10pt,letterpaper]{article}
\usepackage[top=0.85in,left=2.75in,footskip=0.75in]{geometry}

% amsmath and amssymb packages, useful for mathematical formulas and symbols
\usepackage{amsmath,amssymb}

% Use adjustwidth environment to exceed column width (see example table in text)
\usepackage{changepage}

% textcomp package and marvosym package for additional characters
\usepackage{textcomp,marvosym}

% cite package, to clean up citations in the main text. Do not remove.
\usepackage{cite}

% Use nameref to cite supporting information files (see Supporting Information section for more info)
\usepackage{nameref,hyperref}

% line numbers
\usepackage[right]{lineno}

% ligatures disabled
\usepackage[nopatch=eqnum]{microtype}
\DisableLigatures[f]{encoding = *, family = * }

% color can be used to apply background shading to table cells only
\usepackage[table]{xcolor}

% array package and thick rules for tables
\usepackage{array}

% create "+" rule type for thick vertical lines
\newcolumntype{+}{!{\vrule width 2pt}}

% create \thickcline for thick horizontal lines of variable length
\newlength\savedwidth
\newcommand\thickcline[1]{%
  \noalign{\global\savedwidth\arrayrulewidth\global\arrayrulewidth 2pt}%
  \cline{#1}%
  \noalign{\vskip\arrayrulewidth}%
  \noalign{\global\arrayrulewidth\savedwidth}%
}

% \thickhline command for thick horizontal lines that span the table
\newcommand\thickhline{\noalign{\global\savedwidth\arrayrulewidth\global\arrayrulewidth 2pt}%
\hline
\noalign{\global\arrayrulewidth\savedwidth}}


% Remove comment for double spacing
%\usepackage{setspace} 
%\doublespacing

% Text layout
\raggedright
\setlength{\parindent}{0.5cm}
\textwidth 5.25in
\textheight 8.75in

% Bold the 'Figure #' in the caption and separate it from the title/caption with a period
% Captions will be left justified
\usepackage[aboveskip=1pt,labelfont=bf,labelsep=period,justification=raggedright,singlelinecheck=off]{caption}
\renewcommand{\figurename}{Fig}

% Please use the included `plos2025.bst` as your BibTeX style
\bibliographystyle{plos2025}

% Remove brackets from numbering in List of References
\makeatletter
\renewcommand{\@biblabel}[1]{\quad#1.}
\makeatother



% Header and Footer with logo
\usepackage{lastpage,fancyhdr,graphicx}
\usepackage{epstopdf}
%\pagestyle{myheadings}
\pagestyle{fancy}
\fancyhf{}
%\setlength{\headheight}{27.023pt}
%\lhead{\includegraphics[width=2.0in]{PLOS-submission.eps}}
\rfoot{\thepage/\pageref{LastPage}}
\renewcommand{\headrulewidth}{0pt}
\renewcommand{\footrule}{\hrule height 2pt \vspace{2mm}}
\fancyheadoffset[L]{2.25in}
\fancyfootoffset[L]{2.25in}
\lfoot{\today}



%% ---- packages added for this manuscript (permitted: the template forbids
%% only the REMOVAL of its own packages) ----
\usepackage{booktabs}
\usepackage{longtable}
\usepackage{float}
\usepackage{enumitem}
\usepackage{xspace}
\usepackage{etoolbox}

%% PLOS requires double spacing; keep tables single-spaced so they stay legible
\usepackage{setspace}
\doublespacing
\AtBeginEnvironment{tabular}{\singlespacing}
\AtBeginEnvironment{longtable}{\singlespacing}

%% Macros
\newcommand{\maxtoki}{MaxToki-217M\xspace}
\newcommand{\maxtokibig}{MaxToki-1B\xspace}
\newcommand{\geneformer}{Geneformer V2-316M\xspace}
\newcommand{\scfm}{SCFM\xspace}

\newcommand{\PENDING}[1]{\textcolor{red}{\textbf{[PENDING: #1]}}}

%% END MACROS SECTION



\begin{document}
\vspace*{0.2in}

\begin{flushleft}
{\Large
\textbf{Checking agent-run interpretability of a single-cell foundation model: a controlled evaluation of method specifications and a reviewer-derived audit checklist, with a MaxToki-217M case study}
}
\newline
\\
Ihor Kendiukhov\textsuperscript{1*}
\\
\bigskip
\textbf{1} Institute of Medical Genetics and Applied Genomics, University of T\"ubingen, T\"ubingen, Germany
\\
\bigskip

* ihor.kendiukhov@student.uni-tuebingen.de

\end{flushleft}

\section*{Abstract}
We evaluated one way to check interpretability analyses that
large-language-model agents run on biological foundation models. Each
method was written as an agent specification; outputs were audited with
a ten-item checklist of peer reviewers' objections to earlier studies.
An agent used both to run eight interpretability pipelines on the
single-cell foundation model MaxToki-217M and to audit them. We confirmed 149 errors in the deployment's code and outputs, 13
of them critical, including intervention code that deleted a transformer
block and model inputs encoded on the wrong scale. The agent's own audit report found 10\% of a fixed set of 68
errors that predated it, and no critical one; fresh agents in a blind re-audit found 41--45\%, with or
without the checklist. In a controlled study of four known-answer tasks,
all 40 Claude Opus~5.5 runs reached the correct conclusion, with or
without the specification. Executors handled 94.5\% of pitfalls; review and repair added 2.0--2.8 points; generic reviews
turned 3 of 40 correct conclusions wrong. In 11 held-out analyses from
other projects, each in a flawed and a corrected version, every review
of a flawed version found its error, but about 80\% of
reviews of corrected versions claimed a serious error that was not one.
With Opus~5.5, no study showed the checklist finding more errors or
raising fewer false alarms than an equal-budget generic review. Of the 13 critical errors, author-directed agent checks found 8 first and
blind re-audits 3. After correction, MaxToki-217M's attention scores,
sparse-autoencoder features and traced circuits did not beat gene-level
baselines; attention matched co-expression and beat rewired TRRUST
networks in RPE1 cells, not K562. Its gene embeddings aligned with
scGPT's and more with Geneformer's. Hidden states showed no
developmental order beyond cell type; 14 of 15 features steered a maturity score no more than random features.

\clearpage
\newgeometry{top=0.85in,left=1in,right=1in,footskip=0.75in}
\fancyheadoffset[L]{0pt}
\fancyfootoffset[L]{0pt}
\linenumbers

\section*{Introduction}

Single-cell foundation models (\scfm{}s) such as
scGPT~\cite{cui2024scgpt}, Geneformer~\cite{theodoris2023geneformer},
scFoundation~\cite{hao2024scfoundation} and
MaxToki~\cite{maxtoki2026} are trained on millions of single-cell
transcriptomes and are reused for cell-type annotation, gene-network
inference and in-silico perturbation. Whether their internal
representations reflect regulatory biology, or mainly re-organise the
co-expression already present in the training data, cannot be settled by
benchmark accuracy
alone~\cite{kedzierska2023assessing,liu2023evaluating,wenteler2024perteval}.
Mechanistic interpretability asks this question directly: it studies what
a trained model's internal components represent and how they combine,
using probes, sparse autoencoders and causal
interventions~\cite{bricken2023monosemanticity,conmy2023acdc,wang2023ioi,sharkey2025open}.

Such analyses are long chains of code, statistics and interpretation, and
each link can fail. A null model can fail to randomise; an intervention
can edit the wrong tensor; a summary can describe a number as something
it is not. Large-language-model (LLM) agents can now carry out chains of
this length with little human
input~\cite{lu2024aiscientist,boiko2023autonomous,mbran2024chemcrow,yang2024sweagent,chan2024mlebench},
which makes the question of how to check their work practical. Agent
benchmarks for data-driven science and for reproducing published
analyses~\cite{chen2024scienceagentbench,siegel2024corebench,starace2025paperbench}
show that agents often produce plausible results that do not hold up on
inspection. For interpretability of biological models, a wrong result can
be confidently stated, internally consistent and hard for a reader to
detect.

This paper evaluates one answer to that problem. The framework has two
parts. First, each interpretability method is written as a markdown
\emph{specification} that an agent follows when running the analysis.
Each specification is written to a ten-section template, with methods,
parameters, validation criteria and known pitfalls. Seven of the eight
deployed specifications lacked the template's code-references section.
Second, the agent's outputs are checked with an \emph{audit checklist}
of ten methodological objections that peer reviewers raised repeatedly
on the author's earlier interpretability studies. We
deployed the framework on MaxToki-217M, a recently released
trajectory-aware \scfm{}~\cite{maxtoki2026}. An agent ran eight
interpretability pipelines over several working sessions. On the last
day, one of these sessions audited the outputs of all eight with the
checklist.

We then asked how reliable the resulting outputs were, and what the two
parts of the framework contribute. We did this in four ways.
(1)~\emph{Independent verification}: every number in the run summaries was searched for in
the stored run outputs by code, key statistics were re-derived, and the
intervention code was unit-tested; every error found was recorded in an
error ledger with the evidence that confirms it and with who or what found
it first. (2)~A \emph{blind re-audit} (Study~C): fresh agents audited a
copy of the project as it stood before its audit, under three
instructions (the checklist as deployed, a generic version of the
checklist, or a generic critical review of equal budget), and their
findings were scored against the ledger. (3)~A \emph{controlled
execution study} (Study~A): agents performed four analysis tasks built
from the deployment's data, with or without the specification. Each
output was kept as it was, and was also given a generic review and,
separately, a checklist review. After each review, a fresh agent repaired
the output, with the same repair instruction in both review arms. Graders
compared the final conclusions with fixed answer keys. They did not know
the review arm, but they could tell whether the agent had the
specification. (4)~A \emph{held-out benchmark} (Study~B): reviewers with or
without the checklist examined analyses from other projects and other
model families, each in a version with one documented error and in a
corrected version, so that missed errors and false alarms could both be
measured. Finally, we re-ran the main MaxToki-217M analyses with corrected code
and inputs and report what holds for that model. Some deployed analyses,
such as the topology screen and the ageing probe, were not re-run and are
not reported (Results).

\subsubsection*{What this study adds, and what it inherits}
The interpretability methods were developed in the author's earlier
studies of scGPT and
Geneformer~\cite{anon2026attention,anon2026spectral,anon2026topology,anon2026sae,anon2026circuits,anon2026manifold,anon2026exhaustive,anon2026ageing};
MaxToki-217M, the datasets and the agent models are third-party
resources. What was developed for this study is the framework itself (the
specifications and the checklist), its deployment on MaxToki-217M, the
verification and error ledger, the deterministic checks, the three
controlled studies and the corrected MaxToki-217M analyses
(Table~\ref{tab:inherited}). The three studies evaluate the framework;
the MaxToki-217M results are a case study and are stated for that model
only.

{\footnotesize
\setlength{\tabcolsep}{4pt}
\renewcommand{\arraystretch}{1.15}
\begin{longtable}{@{}p{0.30\linewidth}p{0.17\linewidth}p{0.48\linewidth}@{}}
\caption{\textbf{What this study developed and what it inherited.}}
\label{tab:inherited}\\
\toprule
\textbf{Component} & \textbf{Origin} & \textbf{Role in this paper} \\
\midrule
\endfirsthead
\multicolumn{3}{@{}l}{\footnotesize\emph{Table \thetable\ continued from previous page.}}\\
\toprule
\textbf{Component} & \textbf{Origin} & \textbf{Role in this paper} \\
\midrule
\endhead
\midrule
\multicolumn{3}{r@{}}{\footnotesize\emph{continued on next page}}\\
\endfoot
\bottomrule
\endlastfoot
Eight interpretability methods (attention analysis, residual-stream
spectral geometry, topological screening, sparse-autoencoder atlas,
circuit tracing, exhaustive mapping and steering, manifold discovery,
donor-aware ageing probe) & Author's earlier studies of scGPT and
Geneformer & Inherited. The source papers of four of them (attention
analysis, topological screening, sparse-autoencoder atlas, circuit
tracing) are given to agents in Study~A as the source method. \\
MaxToki-217M and MaxToki-1B & Third party~\cite{maxtoki2026} & Deployment
target. \\
Perturb-seq screens, Tabula Sapiens, the Human Lung Cell Atlas, TRRUST, DoRothEA, STRING, GO, KEGG,
Reactome & Third party & Data and reference sets. \\
Claude Code and Claude models & Third party (Anthropic) & Agents in the
deployment and subjects in Studies~A--C. \\
\addlinespace
Pipeline specifications (eight markdown files) & This study (written
at the start of the deployment by Claude Opus~4.7) & First part of the framework; the treatment in
Study~A's contract arm, where the executor also received the deployed
specification. \\
Ten-pattern audit checklist & This study, distilled from six earlier
peer-review comment sets & Second part of the framework; the treatment in
the checklist arms of Studies~A--C. \\
Deployment runs on MaxToki-217M and the in-session audit & This study &
Source of the error ledger and of the Study~C snapshot. \\
Independent verification, error ledger, deterministic checker & This
study & Evaluation of the deployment. \\
Studies~A, B and C (tasks, items, keys, protocol) & This study &
Controlled evaluation of the framework. \\
Corrected MaxToki-217M analyses (fixed intervention hooks, corrected
input encoding, retrained sparse autoencoders) & This study & Case study
of MaxToki-217M. \\
\end{longtable}
}

\section*{Materials and methods}

\subsection*{Study overview}
Fig~\ref{fig:overview} shows how the parts of the study fit together. An
agent ran eight interpretability pipelines on MaxToki-217M by following
the specifications, over several working sessions. Then, in a single
session on the last day, it audited its outputs with the checklist (the
\emph{deployment} and the \emph{in-session audit}). Everything after that was done separately and is the evaluation:
the independent verification that built the error ledger, the three
controlled studies, and the corrected MaxToki-217M analyses. The
controlled studies used fresh agents that had no access to the ledger,
the answer keys, this manuscript or the project's working notes.

\begin{figure}[H]
\caption{\textbf{Study overview.} Left: what was deployed. Each
interpretability method is written as a markdown specification; an agent
follows the specification to run the analysis on MaxToki-217M over
several working sessions and writes run summaries. Then, in one session
on the last day, the agent applies the ten-pattern checklist to the
outputs of all eight pipelines and repairs what it finds. Right: how the deployment was
evaluated. Independent verification (number tracing by code,
re-derivation of statistics from stored per-record outputs, unit tests of
intervention code, an input-encoding audit) and blind re-audits produced
an error ledger of confirmed errors in the deployment's outputs.
Study~C scores blind re-audits of the pre-audit project against that
ledger; Study~A tests the specification and the checklist on four
analysis tasks with known answers; Study~B tests the checklist on held-out
analyses from other projects. The corrected MaxToki-217M analyses re-run
most of the affected pipelines with corrected code and inputs.}
\label{fig:overview}
\end{figure}

\subsection*{The framework}

\subsubsection*{Pipeline specifications}
Each method is described in one markdown file. The template asks for
ten sections, in this order: overview; source (citation and pinned code
commit); inputs; outputs; dependencies; methodology (numbered steps with
equations and hyperparameters); code references (pointers into the pinned
reference code); a parameter table with defaults and ranges; validation
criteria; and known pitfalls. Only one of the eight deployed
specifications (the ageing-probe one) has all ten. The other seven have
no separate code-references section. Their code pointers are mostly inside
the methodology section, and they end with a quick-start section on
applying the method to a new model (see \emph{What deterministic checks
catch} in Results). The manifold specification names no pinned commit,
because its source paper links no code. The specifications for the eight
pipelines were written on 15 April 2026, at the start of the deployment,
by Claude Opus~4.7 under the author's direction, from the source papers
and their code; they are 48,000--96,000 characters long. During the deployment the ten-section
structure was an instruction to the agent; no code read or checked the
specifications. A specification therefore worked as a long, structured
prompt that the agent was asked to follow. The executing agent chose, for
example, the number of null-model draws and the number of cells where
compute was limited, and recorded some but not all of these choices.

\subsubsection*{Per-model project folders}
Model-specific code (loading the checkpoint, mapping genes to tokens,
encoding cells as model input) lives in a project folder separate from the
specifications, so that a specification can be pointed at a new model
without being rewritten. Each pipeline run writes its scripts, logs,
outputs and an agent-written run summary into its own folder.

\subsubsection*{The ten-pattern audit checklist}
The checklist (Table~\ref{tab:checklist}) was distilled from six
peer-review comment sets on the author's earlier interpretability papers.
For each pattern it gives the objection as a reviewer would state it, a
diagnostic signature to search for in run outputs and code, and a
prescription. The deployed version is a full audit specification.
Besides the ten patterns, it has the other template sections and a
step-by-step audit workflow: list the runs, check each pattern in each
run, compare runs, grade severity, propose actions and write a report.
For each pattern it also lists examples from the MaxToki project and cites
the review documents that raised it. Studies~A and B, and one arm of
Study~C, used a shorter \emph{generic} version, about 40\% of the length
of the deployed file. It keeps the ten pattern entries, without the
project-specific examples and the citations, and with a few
project-specific phrases made general. It adds a short opening
instruction and leaves out the workflow and the other sections. Another
arm of Study~C received the deployed file. Both versions are in the code
repository.

{\footnotesize
\setlength{\tabcolsep}{4pt}
\renewcommand{\arraystretch}{1.15}
\begin{longtable}{@{}p{0.17\linewidth}p{0.26\linewidth}p{0.28\linewidth}p{0.24\linewidth}@{}}
\caption{\textbf{The ten-pattern audit checklist.}}
\label{tab:checklist}\\
\toprule
\textbf{Pattern} & \textbf{Objection} & \textbf{Signature searched for} &
\textbf{Prescribed repair} \\
\midrule
\endfirsthead
\multicolumn{4}{@{}l}{\footnotesize\emph{Table \thetable\ continued from previous page.}}\\
\toprule
\textbf{Pattern} & \textbf{Objection} & \textbf{Signature searched for} &
\textbf{Prescribed repair} \\
\midrule
\endhead
\midrule
\multicolumn{4}{r@{}}{\footnotesize\emph{continued on next page}}\\
\endfoot
\bottomrule
\endlastfoot
\textbf{P1} Weak null or no chance anchor & ``What would chance alone give?''
& A number without its chance value; a null that only shuffles labels; a
$p$-value without a named null. & Report the null mean and spread next to
every load-bearing number; prefer the strictest null. \\
\textbf{P2} No trivial baseline & ``What does a simpler method give?'' &
A performance figure without a comparison to expression-only or
gene-level baselines. & State every capability as a gain over the
simplest comparator. \\
\textbf{P3} Endpoint does not match the claim & ``The metric measures
something else.'' & A regulatory claim resting on a ranking score; the
same database used to select and to evaluate. & State what the endpoint
literally measures; weaken the claim to match. \\
\textbf{P4} Wrong unit of inference & ``Your unit is donors, not
cells.'' & $p$-values over cells for donor-level claims; random
cross-validation splits when rows share a factor. & Name the unit;
group resampling and splits by the shared factor. \\
\textbf{P5} Selection-driven inflation & ``You selected on one property
and tested another.'' & Hand-picked subsets; ranking and testing on the
same data. & Separate selection from evaluation; control the selected
property. \\
\textbf{P6} No positive control & ``Could the pipeline detect a signal at
all?'' & A negative result with no planted or known-true signal. & Run a
synthetic and a real positive control first. \\
\textbf{P7} Generalisation from one condition & ``One cell line, one
architecture.'' & Conclusions from one dataset, tissue or model without
a scope statement. & Replicate on a held-out condition or state the
scope. \\
\textbf{P8} Unknown stability & ``What is the interval? Another seed?'' &
Point estimates without intervals, threshold sweeps or seed checks. &
Bootstrap at the unit of inference; report threshold and seed
sensitivity. \\
\textbf{P9} Causal overreach & ``Say `is consistent with', not
`confirms'.'' & Causal verbs on correlational evidence. & Weaken the verb
or run the intervention. \\
\textbf{P10} Missing reproducibility scaffolding & ``Provide code,
environment and compute.'' & Cited scripts absent; no pinned environment,
seeds or compute record. & Freeze environment and seeds; publish every
script; report compute. \\
\end{longtable}
\begin{flushleft}
Each row is an
objection that reviewers raised repeatedly on the author's earlier
interpretability papers, rewritten so that an agent can search for it.
AUROC is the area under the receiver-operating-characteristic curve (0.5
is chance); a null model is a randomised version of the data that keeps
some structure and destroys the effect under test.
\end{flushleft}
}


\subsubsection*{What was checked by code, by agent judgment, or by a person}
Table~\ref{tab:checks} lists each requirement the framework imposes and
how it was checked during the deployment. Apart from pass/fail flags that
two pipelines computed from their own numbers, every check was an agent
judgment; the author's role was to start runs, answer questions, and ask
whether issues had been addressed. To make the distinction testable we
wrote a deterministic checker (see \emph{Deterministic checks} below) and
report what it can and cannot detect.

{\footnotesize
\setlength{\tabcolsep}{4pt}
\renewcommand{\arraystretch}{1.15}
\begin{longtable}{@{}p{0.22\linewidth}p{0.36\linewidth}p{0.36\linewidth}@{}}
\caption{\textbf{How the framework's requirements were checked.}}
\label{tab:checks}\\
\toprule
\textbf{Requirement} & \textbf{During the deployment} & \textbf{Deterministic checker} \\
\midrule
\endfirsthead
\multicolumn{3}{@{}l}{\footnotesize\emph{Table \thetable\ continued from previous page.}}\\
\toprule
\textbf{Requirement} & \textbf{During the deployment} & \textbf{Deterministic checker} \\
\midrule
\endhead
\midrule
\multicolumn{3}{r@{}}{\footnotesize\emph{continued on next page}}\\
\endfoot
\bottomrule
\endlastfoot
Specification has the ten sections, in order & Instruction to the agent;
not checked & Detects missing or misordered sections. \\
Pinned code and code references resolve & Agent wrote commit hashes; not
checked & Detects unresolvable pins, malformed references, missing paths
and lines. \\
Parameters used equal the specification, or the deviation is justified &
Agent's choice; partly noted in code comments & Detects deviations from
machine-readable defaults and ranges where the run records the value;
cannot judge whether a deviation is acceptable. \\
Null models actually randomise & Agent judgment (audit pattern P1) &
Checks null size only; cannot tell whether a null randomises. \\
Trivial baselines and positive controls exist & Agent judgment (P2, P6) &
Detects whether the specification has a labelled slot; cannot judge
adequacy. \\
Intervals name their method and resampling unit & Agent judgment (P4, P8)
& Points to interval statements without a method; cannot judge
correctness. \\
Wording matches the evidence & Agent judgment (P3, P9) & Lists causal
words for a reader; cannot judge context. \\
Every reported number comes from the run's outputs & Not checked & Searches the stored outputs for
each number and gives the rate of chance matches; cannot judge
meaning. \\
Pre-registered quality gates are computed and not loosened & Two
pipelines computed their own pass/fail flags & Recomputes stored gate
flags and detects uncomputed gates. \\
Seeds, environment, device and checkpoint are recorded & Agent judgment
(P10) & Detects scripts that use random numbers without a seed, and runs
with no device, package-version, checkpoint or run-manifest record;
cannot tell whether every random call uses the seed. \\
Hypothesis retirement in automated loops & Described in the
specifications (retire after two negatives in a row); not coded in any
pipeline. A scripted batch of six hypotheses in the topological screen
coded its own promote, inconclusive or retire label, with a stop after
three retires in a row that was never reached; a manual review later
changed 4 of the 6 labels & Detects whether a decision log exists. \\
Go/no-go decisions on expensive runs & Author's prompts; not logged &
Not checkable from files. \\
\end{longtable}
\begin{flushleft}
``During
the deployment'' describes what was in place when the eight runs and the
in-session audit were carried out. ``Deterministic checker'' describes
what the checker written afterwards can detect by reading files; what it
cannot detect remains a matter of judgment.
\end{flushleft}
}


\subsection*{The MaxToki-217M deployment}

\subsubsection*{Model and data}
MaxToki-217M~\cite{maxtoki2026} is a Llama-architecture autoregressive
transformer with 11 transformer blocks, hidden size 1,232, 8 attention
heads per block and a vocabulary of 20,275 tokens, 20,271 of them genes (Hugging Face
\texttt{theodoris-lab/MaxToki}). The downloaded weight files match
revision \texttt{21aa7b7f} byte for byte. That was the \texttt{main}
branch on the download day. The same bytes are also at earlier and later
revisions, so the hashes confirm the weights but not the exact revision.
The model exposes 12 hidden states, called layers 0--11 below. Layer~0
is the token embedding. For $l$ from 0 to 10, layer~$l$ is the input of
transformer block~$l$ (blocks counted from 0). Layer~11 is the output of
the last block after the final normalisation. Layer-8 attention is the
attention inside block~8. One attention run also used MaxToki-1B. Every pipeline gave
the model one cell per sequence (about 1,000--2,900 tokens per cell); no
multi-cell trajectory inputs were used. Cells came from the Replogle
essential-gene Perturb-seq screens in K562 and
RPE1~\cite{replogle2022perturbseq} (the K562 cells from the filtered public
copy in the Hugging Face dataset \texttt{arcinstitute/State-Replogle-Filtered},
and the RPE1 cells from the scPerturb copy~\cite{peidli2024scperturb}),
the Adamson CRISPR-interference
screen~\cite{adamson2016perturbseq} (in the processed form distributed with
GEARS~\cite{roohani2024gears}), Tabula
Sapiens~\cite{tabulasapiens2022} and, for the external-lung panel of the
topological hypothesis screen only, 1,500 cells from the Krasnow Lab Human
Lung Cell Atlas~\cite{travaglini2020lung} (Smart-seq2 data; CELLxGENE dataset
version \texttt{c88e0403-da93-40f4-99b5-f5fdeb81a82c}). Reference sets were
TRRUST~v2~\cite{han2018trrust}, ChIP-seq target sets from
DoRothEA~\cite{garcia-alonso2019dorothea}, the STRING protein-interaction
network~\cite{szklarczyk2023string}, and GO
biological-process~\cite{ashburner2000go}, KEGG~\cite{kanehisa2000kegg}
and Reactome~\cite{gillespie2022reactome} gene sets. In the
deployment, DoRothEA was used only in the audit's repairs of the
transcription-factor test; the corrected analyses use it for the same
test. STRING was used by the spectral-geometry, topology and SAE-atlas
runs. The gene sets were used mainly to annotate SAE features. The GO,
KEGG and Reactome sets are the Enrichr libraries
GO\_Biological\_Process\_2023, KEGG\_2021\_Human and Reactome\_2022, as
named in the script of an earlier project that built them. The
spectral-geometry and SAE-atlas runs also read GO terms, including
cellular-component terms, from a gene-to-GO table that appears to be the
one distributed with GEARS~\cite{roohani2024gears}. No file records which
release of DoRothEA or STRING, or which version of that gene-to-GO table,
was used.

\subsubsection*{Pipelines, compute and supervision}
Table~\ref{tab:pipelines} lists the eight pipelines and what was run.
All model computation of the deployment and of the corrected analyses ran
on one Apple-silicon MacBook Pro laptop with
32~GB unified memory. It used the PyTorch Metal (MPS) backend in 32-bit
floating point; some analysis ran on the same machine's CPU. The run
files do not name the chip. The machine that holds the files today has
an Apple M2 Pro and is probably the same one, but no file proves it.
From the first to the last run file the deployment spanned 21.9 days (15 April to 7 May 2026), with
files written on 11 of those days. Log files cover 47.8~hours of
computation when overlapping runs are counted once. The agent was Claude
Opus~4.7, run through Claude Code (Anthropic). The run files do not
record the model, and the session transcripts were not kept; the model is
known from the author, and a usage note that the agent wrote on 8 May
2026 also names it. The one automated hypothesis loop, in the
spectral-geometry pipeline, called Claude Sonnet (version not recorded).
The author typed 73
prompts during the deployment. Supervision time was not logged.

{\footnotesize
\setlength{\tabcolsep}{4pt}
\renewcommand{\arraystretch}{1.15}
\begin{longtable}{@{}p{0.20\linewidth}p{0.28\linewidth}p{0.36\linewidth}r@{}}
\caption{\textbf{The eight pipelines as run on MaxToki-217M.}}
\label{tab:pipelines}\\
\toprule
\textbf{Pipeline} & \textbf{Question} & \textbf{What was run} & \textbf{Log h} \\
\midrule
\endfirsthead
\multicolumn{4}{@{}l}{\footnotesize\emph{Table \thetable\ continued from previous page.}}\\
\toprule
\textbf{Pipeline} & \textbf{Question} & \textbf{What was run} & \textbf{Log h} \\
\midrule
\endhead
\midrule
\multicolumn{4}{r@{}}{\footnotesize\emph{continued on next page}}\\
\endfoot
\bottomrule
\endlastfoot
Residual-stream spectral geometry & How does each layer's representation
of genes change with depth? & Per-layer singular-value geometry,
effective rank and stability across cell samples; 2,000 TS immune cells;
a later cross-model comparison with Geneformer V2-316M. & 3.6 \\
Attention-derived regulatory networks & Do attention maps recover
regulator--target pairs? & Layer-8 attention against TRRUST and against
knockdown responses; four runs (K562, RPE1, Adamson with MaxToki-217M;
K562 with MaxToki-1B). & 8.8 \\
Topological hypothesis screen & Does higher-order geometry of gene pairs
carry regulatory information? & A subset of the catalogue (21 of 141
hypotheses had code); the strictest null at 3 of 12 layer states; a
cross-model comparison with scGPT. & 2.6 \\
SAE atlas & What do sparse features of the residual stream represent? &
TopK SAEs (4,928 features, $k=32$) at all 12 layer states, trained on 500
K562 cells; annotation; a transcription-factor specificity test;
feature patching. & 10.1 \\
SAE circuit tracing & Which feature-to-feature edges carry causal
influence? & Ablation-based edges from 30 features at each of four layers
(200 K562 cells); a CRISPRi direction-of-change test. & 4.2 \\
Exhaustive mapping and steering & Is the computation built from
combinatorial logic, and can features steer cell state? & Edges from
1,000 layer-5 features (20 cells); four feature triplets (50 cells);
steering at five layers (50 TS immune cells). & 7.7 \\
Manifold discovery & Is there a recoverable developmental ordering of
cells? & A pre-set hematopoietic ordering and 18 candidate orderings,
scored with quality gates on cell-group centroids from TS donors. & 15.2 \\
Donor-aware ageing probe & Is donor age encoded? & First stage only; its
gate failed. & 1.5 \\
\end{longtable}
\begin{flushleft}
``Log
hours'' is the span covered by each pipeline's log files, with
overlapping runs counted once; it is laptop wall-clock time, not GPU
time, and can include idle time. Model computation ran on one laptop;
the runs that recorded a device used the Apple-silicon MPS backend (one
also used the CPU). The 12 layer states are the hidden states the model
exposes: the token embedding (layer~0) and the outputs of its 11
transformer blocks (layers 1--11; layer~11 is taken after the final
normalisation). TS, Tabula Sapiens; SAE, sparse autoencoder.
\end{flushleft}
}


\subsubsection*{The in-session audit}
On the last day of the deployment the agent wrote the audit specification
and applied the checklist to all eight pipelines in one working session.
Earlier that day the same session had extended the topological hypothesis
screen, so for that pipeline it graded its own same-day work. All other
runs, including the first run of that screen, had been made in other
sessions. The agent wrote an 8~$\times$~10
matrix of free-text verdicts (pipeline $\times$ pattern) and a list of
proposed repairs, and then carried out two rounds of repair, each started
by a prompt from the author. From the audit request to the last repair
took 3.7~hours. The audit sweep read mainly the run summaries, plus
selected output tables and scripts, and re-ran no model. Three repair
steps did run MaxToki-217M again: a re-extraction of the lung cells with a
new sampling seed, a targeted re-run of the transcription-factor test for
GATA1, and a small cell-level benchmark.

\subsection*{Independent verification and the error ledger}
After the deployment, each claim in the run summaries and each number
they report was checked in four ways. (i)~A deterministic number tracer
searched each run's stored outputs for every number in its summary,
together with a chance rate for such matches. (ii)~Key statistics were
re-derived from stored per-record outputs with new code (for example,
direction-of-change predictions were rescored against constant-sign
baselines on the identical records; null models were re-sampled and
checked for actual randomisation). (iii)~Intervention code was read
against the transformer library source and unit-tested: an edit of zero
must leave the output unchanged, an edit at layer~$l$ must leave earlier
layers unchanged, and combined edits must differ from single edits.
(iv)~The model-input encoding of every run was compared with the
encoding the model expects. Claude Opus~5.5 agents ran checks (i)--(iv)
in working sessions that the author started and directed. Some of these
checks followed up points raised by a human reader from outside the
project (not the author). So an error credited to independent verification was found
by an agent doing a directed check, not by an agent working on its own.
Further errors came from the blind re-audits
of Study~C: every finding that matched no ledger entry was checked against
the full project by a separate agent, which confirmed or rejected it with
code. This agent's confirming code was not kept; the ledger rows give
the recomputed values. Four errors were first pointed out by the human reader.
Independent verification then confirmed each of them from the stored
files and code.

Each confirmed error is one row of the ledger (S1~Table), with its
location, the deterministic evidence that confirms it, its type (code
bug, wrong statistic, missing baseline or null, wrong description,
unsupported claim, reproducibility), the checklist pattern it falls
under (or ``outside the checklist''), whether it was present before the
in-session audit (from file timestamps), who or what first found it, and
whether the in-session audit report caught it. The full ledger in the code
release also records who or what confirmed each error. Severity was graded as \emph{critical} (a headline result or claim cannot
stand), \emph{major} (it stands only with material qualification) or
\emph{minor}. The ledger covers the deployment's code, logs, outputs, run
summaries and audit reports.

\subsection*{Subject agents and context control}
All agents in the pre-registered parts of the three controlled studies
(executors, reviewers, repairers and graders) were Claude Opus~5.5
(\texttt{claude-opus-5-5}). In the exploratory arm of Study~A (described
below), the executors, reviewers and repairers were Claude Haiku~4.5 and
the graders were Opus~5.5. All agents were run as Claude Code workflow
subagents of the built-in read-only type with the model set explicitly.
We used this agent type because it starts
without project instruction files and without the author's memory notes,
both of which contain descriptions of earlier errors; we checked this with
a probe agent before the studies. The agents can read files and run
shell commands. They were asked to run all analyses as inline Python
(\texttt{numpy}, \texttt{pandas}, \texttt{scipy}, \texttt{scikit-learn};
CPU only, at most 9~minutes per command) without writing files, and to
read only inside their task folder. Deliverables were returned as
structured output and written to disk by the harness. Study~A task folders and Study~B item folders were placed
at a neutral path under opaque names, so that folder names revealed
neither the study, the arm nor the status of an item. After each study we
scanned every tool call of the subject agents in its pre-registered part:
341 executors, reviewers, repairers and re-auditors in all (200 in
Study~A, 132 in Study~B and 9 in Study~C). Graders were meant to read the
keys and were not scanned. We looked for any access to the project
repository, the answer keys, the protocol files other than the checklist
given to a reviewer, the author's notes, this manuscript, or the folder
of another task or item. No subject agent accessed any of them. The
subject agents of the exploratory Haiku arm were checked in a separate
scan (see below and S2~Appendix). We also checked whether any agent
read a file that another agent had written. Twenty-five agents used a
scratch folder that the session shares between agents. Three of them
wrote temporary files there. The other 22 only ran a command to make sure
the folder existed. Six of these 22 also listed the folder and so saw the
names of other agents' files; none opened such a file.

In Study~A task T1 (whether layer-8 attention encodes TRRUST
regulator--target pairs in RPE1 cells), a Claude
Haiku~4.5 executor from the first run of the exploratory arm left a
script in the task folder. That run was set aside (see \emph{Study A}
below). All ten Opus~5.5 repair agents of T1 in the contract arm (the
agents that revised a deliverable after its review) opened this script. The contract arm is the Study~A arm in which the
executor also received the deployed specification. All ten of their
deliverables reached the correct verdict. Leaving out the T1
contract-arm runs does not change any Study~A result. For example, a
checklist review followed by repair still raised the share of pitfalls
handled by 2.9 percentage points over no review (95\% bootstrap interval
over the remaining 35 executor runs, 0.7 to 5.1). Graders of these
deliverables also opened that script.
Every task and item folder was afterwards checked to be identical to its
source.
Reasoning effort was the same session default for all agents.

The design, outcomes and analyses were written down before the studies
were run, and the answer keys were frozen by hash before the studies that
used them (S2~Appendix). We call this protocol pre-registered in that
sense: it was frozen by file hash and logged before the studies, but it
was not lodged in a public registry. The protocol has three amendments, all given in
S2~Appendix. Amendment~1 was logged on 1 October 2026, after Study~C had run
and before Studies~A and B started. It replaced one Study~A task, split
Study~A into two batches with the same protocol, and excluded one Study~B
pair that failed its own verification. Amendment~2 added the exploratory
Haiku arm described below, after the results of Study~A task T2 were
known. Amendment~3, logged before that arm was re-run, set aside the
arm's first run and gave every agent its own folder.

\subsection*{Study C: blind re-audit of the deployment}
We used file timestamps to rebuild a snapshot of the project as it stood
before the in-session audit started. The snapshot held all run code, logs
and configuration files, the specifications and the run summaries. It
also held most outputs smaller than 5~MB, and nine outputs of 5--20~MB
chosen because they hold the per-gene-pair or per-cell-group data behind
headline numbers. Every other output was replaced by a listing of its
name, size and shape. Where the audit had edited a summary in place, we
removed its dated audit additions. The audit files, all manuscript folders and
every later file were excluded. The reference set was the 68 ledger
errors that had been logged when the study started and were present in
this snapshot (6 critical, 25 major, 37 minor). When the set was frozen,
60 of them had been confirmed. The other 8 (3 critical, 3 major, 2 minor)
were confirmed later: 6 by new analyses, 1 by reading the code and 1 by
a literature check. All 68 are in S1~Table. Three arms received
the same project description, output format and budget instruction
(``Work as a careful expert would in about one working day. You may use
up to roughly 150 tool calls''): the checklist as deployed (the full
audit specification, with its project-specific examples and its step-by-step audit workflow, which
the prompt asked the agent to follow), the generic checklist (the ten
pattern entries only), or an instruction to review the project critically
before publication. Each arm was run three times.

Two graders, given the ledger and the audit's findings, judged for every
ledger error whether the audit detected it (\emph{yes}: a finding
identifies the same problem clearly enough that acting on it would lead
to the fix; \emph{partial}: it points at the same result or code with a
related but wrong or incomplete concern; \emph{no}), and matched every
finding to ledger rows. A third agent resolved disagreements. Findings
that matched no ledger row were checked by a fourth agent against the
snapshot and classified as real error, real but trivial, false alarm, or
undecidable; real errors were then verified against the full repository
and added to the ledger. The in-session audit was scored on the same 68
errors from its written sweep report, made before the repair rounds. The
Study~C graders did not score it. Its score is the ledger's record of
whether that report caught each error (yes, partly or no). The
independent verification made this record when it built the ledger,
before Study~C ran. The record was not blinded, and no second grader
checked it. ``Partly'' means that the report flagged
the area but not the actual error.

\subsection*{Study A: controlled execution study}
Four analysis tasks were built from the deployment's data
(Table~\ref{tab:tasks}). Each task asks a question about MaxToki-217M
that the deployment had answered, gives the data needed to answer it,
and asks for a report with a verdict, estimates with uncertainty, and a
complete analysis script. Each task has known pitfalls into which the
deployment fell (for example, a 50\% baseline for a direction-of-change
score whose classes are unbalanced). The pitfalls are not mentioned in
the task brief.

{\footnotesize
\setlength{\tabcolsep}{4pt}
\renewcommand{\arraystretch}{1.15}
\begin{longtable}{@{}p{0.06\linewidth}p{0.30\linewidth}p{0.22\linewidth}p{0.36\linewidth}@{}}
\caption{\textbf{Study A tasks.}}
\label{tab:tasks}\\
\toprule
& \textbf{Question} & \textbf{Key verdict} & \textbf{Main pitfalls (not stated in the brief)} \\
\midrule
\endfirsthead
\toprule
& \textbf{Question} & \textbf{Key verdict} & \textbf{Main pitfalls} \\
\midrule
\endhead
\bottomrule
\endlastfoot
T1 & Do layer-8 attention scores encode TRRUST regulator--target pairs
(RPE1 cells)? & Not supported (``inconclusive'' accepted under stated
conditions) & Gene-level properties of targets; a degree-preserving null
that actually rewires; few transcription factors as the resampling unit;
model-free pair signals. \\
T2 & Do SAE-circuit predictions get the direction of change after a
CRISPRi knockdown right (K562)? & Not supported (``inconclusive''
accepted under stated conditions) & Unbalanced classes (50\% is not
chance); constant-sign predictors; records clustered by silenced gene; a
model-free target-direction baseline. \\
T3 & Are SAE features that respond to a transcription factor's knockdown
specific to its target genes? & Not supported & Rare and long genes
dominate feature gene lists; specificity requires comparison with other
factors' targets; threshold search; one factor is not the question. \\
T4 & Is MaxToki-217M's gene-embedding geometry aligned with scGPT's, and
how does that compare with Geneformer V2-316M? & Aligned with scGPT, but
clearly less than with Geneformer & scGPT rows must be looked up by
token id; in-sample canonical correlation has a high chance level; genes
are the resampling unit. \\
\end{longtable}
\begin{flushleft}
Each task's data are a fixed snapshot of
the deployment's outputs, except T2. The deployed circuit edges were made
by an intervention hook that deleted a transformer block, so T2's circuit edges come from a re-run
of the deployed circuit tracing with corrected hooks (S2~Appendix,
Amendment~1). The re-run used the same cells, input encoding, sparse
autoencoders and source features as the deployment. Each answer key
refers to its task's data. T2 and T3 use sparse-autoencoder features trained on the deployment's incorrectly
encoded inputs (see \emph{Errors in the deployment's outputs} in
Results); this does not affect the tasks as analysis exercises, because
their keys were computed on the same data.
\end{flushleft}
}


In the \emph{paper} arm the agent received the brief, the data and the
source paper of the method. In the \emph{contract} arm it also received
the deployed specification and one sentence asking it to follow it (the
arm names are those of the pre-registration, S2~Appendix). Both arms
therefore had the same task information and the same access to the
source method. Each executor output was then (a)~left as is,
(b)~reviewed by a fresh agent instructed to review it critically, or
(c)~reviewed by a fresh agent instructed to apply the generic checklist;
in (b) and (c) a fresh agent then revised the deliverable given the
review points, with the same repair instruction in both arms. Review
prompts had the same budget sentence (``Spend about the effort of a
careful two-hour expert review; you may use up to roughly 60 tool
calls.''). Each task--arm cell had five
independent executor runs, giving 40 executor runs and 120 final
deliverables.

Each answer key states the correct verdict and any conditions under which
a hedged verdict also counts as correct, key numbers with tolerances,
six pitfalls, each with what correct and incorrect handling looks like,
and tempting false statements. Each key was written by one agent. Its key numbers were computed by a
script that reads only the task's data. A second agent re-derived the key
numbers with its own code, and the keys were then frozen by hash. Two graders scored each deliverable against the key;
a third grader decided when the first two disagreed on the primary
outcome. Before grading, words that could reveal the arm (such as
``specification'', ``review'' and ``repair'') were replaced by a
placeholder in the deliverable. Graders also read the task brief, which in
the contract arm points to the specification file, so they could tell the
executor arm. Only the review arm was hidden from them.

The primary outcome was a correct verdict. Secondary outcomes were the
share of pitfalls handled (yes~=~1, partly~=~0.5, no~=~0), the number of
materially false statements, key numbers within tolerance, the share of
review points that a grader judged invalid, verdict changes caused by
repair, and whether the returned analysis script reproduced the reported
estimates. For this last outcome, the harness ran each script on the
task data with the same Python interpreter that the subjects used, with a
20-minute limit. Scripts that stopped while other scripts shared the
machine were run again, alone or two at a time, with a 60-minute limit;
none of these needed more than 10~minutes. A reported estimate counted as
reproduced when the script printed a number of the same name that agreed
with it within 0.005 or within 1\% of its value, whichever was larger.

There were four primary contrasts. The first compared the contract and
paper arms without review. It used a Fisher exact test, the risk
difference with a Newcombe interval~\cite{newcombe1998interval}, and a
Mantel--Haenszel test stratified by task. The other three compared review
arms on the same executor outputs: checklist versus generic, checklist
versus none, and generic versus none. These used exact McNemar
tests~\cite{mcnemar1947note}. We applied Holm
correction~\cite{holm1979simple} across all four contrasts.

Because every Opus~5.5 deliverable reached the correct verdict on the
first task to finish (T2), we added an exploratory arm after that task's
results were known, while the other tasks were still running
(Amendment~2): the same design with Claude Haiku~4.5 as executor, reviewer
and repairer, and the same Opus~5.5 graders and keys. A first run of this
arm was set aside and not analysed. Its Haiku agents did not keep to the
instruction not to write files: they wrote scripts into the shared
temporary folder and into shared task folders, where other agents could
read them, and some did. The arm was then run again with one folder per agent, each holding
links to the task files (Amendment~3, S2~Appendix). Only the re-run is
reported. S2~Appendix describes the checks of which files its agents
read. Five reviewers opened the folder of the executor whose deliverable
they were reviewing, because the deliverable named it. Every read of a
file in the shared temporary folder was of a file that the reading agent
had written itself. The arm is reported separately and does not enter the
pre-registered contrasts.

\subsection*{Study B: held-out audit benchmark}
Study~B used analyses that played no part in building the checklist.
They came from three other projects by the author. One studied
Cell2Sentence (C2S-Scale-Gemma-2-2B)~\cite{rizvi2025c2sscale} cell
embeddings. One studied the ESM-2 (650M)~\cite{lin2023esm2} protein
language model. One studied frozen scGPT, Geneformer V2-316M and
STATE-SE~\cite{adduri2025state} cell embeddings, together with
gene-regulatory-network and perturbation benchmarks, including the
Perturb-seq data of Norman et al.~\cite{norman2019exploring} as
distributed by scPerturb~\cite{peidli2024scperturb}. One pair used the
fetal human gut atlas of Elmentaite et al.~\cite{elmentaite2020gut}, and
three used proteins and residue annotations from
UniProtKB/Swiss-Prot~\cite{uniprot2025}. Four of the 11 pairs
in the study use no model at all (two label-building steps and two
benchmarks). Each analysis was rebuilt as a self-contained package (data,
code, outputs and a results summary written as the original analyst would
have written it) in a \emph{flawed} version with exactly one documented
error and a \emph{clean} version that differs only in that error. Of
twelve planned pairs, nine had a natural error that had been found and
fixed after the fact. The other three had a planted overclaim of causal
or mechanistic meaning in the summary (the checklist's pattern P9, for
which the pool had no natural example). Every package was checked by
code: the flawed version reproduces its wrong result, the clean version
the correct one, and comments or names that hinted at the error were
removed. One natural pair failed this check and was dropped, leaving 22
items (11 flawed, 11 clean; 6 errors inside the checklist's patterns and
5 outside, such as a circular statistic that breaks on evenly spread
angles or a row-index join that matched the wrong cells).

Each item was reviewed three times with the generic checklist and three
times with a generic critical review, with the same budget sentence in
both arms (``Spend about the effort of a careful one-hour expert review;
you may use up to roughly 50 tool calls.''). Two graders with the item's key judged
whether a review detected the documented error under a pre-written
detection rule, and classified every finding as the documented error,
another real problem, a false alarm, a minor or style point, or disputed;
a third grader resolved disagreements, and a verifier agent resolved
disputed findings by running code on the package. There were three
outcomes. The first was sensitivity: the share of reviews of flawed items
that found the documented error. The second was the share of reviews of
clean items with at least one false alarm. A finding counted as a false
alarm in one of two ways. In the first, the graders judged that it
claimed a problem that does not exist, checked against the item's key and
the data, and that it said a headline result or claim was wrong. In the
second, a grader marked the finding as disputed and the verifier agent,
after running code on the package, ruled it a false alarm. These findings
counted at any severity, whether or not they concerned a headline result.
The third was the number of other real problems found per review.

After the results were known, we also sorted the false alarms about a
main result that the graders had marked into four groups: correct but
changing no stated conclusion; correct and already stated in the
analysis's own summary or README; factually wrong; or a real problem that
the key had missed. This step was not pre-registered. Each claim was
labelled by one agent, with no second rater. False alarms that the
verifier agent ruled after a grader dispute were not sorted.

\subsection*{Corrected MaxToki-217M analyses}
The verification found that the intervention code of two pipelines (SAE
circuit tracing, and exhaustive mapping and steering) edited the wrong
tensor and that every run except the RPE1 attention run encoded the model
input incorrectly (Results). We therefore re-ran the main affected
analyses (those not re-run are listed in Results) with three changes.
First, the model input follows MaxToki's rank-value encoding. Each gene's count is divided by that gene's median in the
Geneformer gene-median table (gc104M) that the MaxToki tokenizer uses,
and genes are ranked from high to low. The K562 and Adamson files store
only $\log(1+\text{CP10k})$ values, where CP10k means counts scaled so
that each cell sums to 10,000. For these files we undid the log with
$\exp(x)-1$, divided by the smallest non-zero value in the cell and
rounded. This gives whole-number counts up to one factor per cell, so the
gene ranking is the same as from the raw counts. For the RPE1, Tabula
Sapiens and lung files we used the stored raw counts. Second, the
intervention hooks edit the input of transformer block~$l$, which is the
tensor that the layer-$l$ sparse autoencoder reads. They add each edit to
the live activations, so edits at several layers add up instead of
overwriting each other. Third, we retrained all twelve sparse
autoencoders (one per layer state) on correctly encoded cells. Every
model input was checked against the expected token order before each
forward pass. Each re-run that edits the model had sanity checks (zero
edit gives zero effect; edits change only later layers; combined edits
differ from single ones; a feature that never fires changes nothing). The
key numbers of each corrected analysis were computed a second way with
separate code. Other agents also re-derived them with their own code:
three for the RPE1 attention analysis and one for the cross-model
comparison (S3~Appendix). Methods
specific to each analysis are given with its result and in S3~Appendix.

\subsection*{Deterministic checks}
The checker has three parts: a specification linter (presence and order
of the ten sections, a resolvable pinned commit, code references in the
required form, a machine-readable parameter table), a run-manifest check
(parameters used versus the specification's defaults and ranges, gate
flags recomputed from stored values) and the number tracer described
above. It reads files only and runs no model. Its code and outputs are in
the repository.

\subsection*{Statistics}
Proportions are reported with Wilson 95\% intervals~\cite{wilson1927probable}
unless stated otherwise. Bootstrap intervals are percentile intervals.
The one exception is the feature-triplet analysis, which uses basic
bootstrap intervals over cells, as stated with its result. The resampling
unit is the item in Study~B, the reference error in Study~C, the executor run in Study~A, the
deployed pipeline run (eight in all) for the number tracer, and the knockdown
(silenced gene), transcription factor, gene, donor or cell in the
MaxToki-217M analyses. The spectral-geometry analysis uses jackknife
intervals instead, over blocks of cells or groups of genes. In Study~A
the main intervals resample executor runs pooled over all four tasks. The
pre-registered plan resamples runs within each task. Results report
intervals of that kind as well, for the main study and for the Haiku~4.5
arm. Study~A bootstrap intervals use 4,000 draws; those of Studies~B and
C use 2,000. In the MaxToki-217M analyses, multiple testing within a
family of tests was controlled with the Benjamini--Hochberg (BH)
procedure~\cite{benjamini1995fdr}. A BH-adjusted $p$-value is called a
$q$ value, and a test passes when $q < 0.05$. Study~A used Holm
correction instead, as described above. Empirical $p$-values use the
$(b+1)/(m+1)$ form~\cite{phipson2010permp}. Degree-preserving network
nulls used the Curveball algorithm~\cite{strona2014curveball} and were
checked to change the network in every draw. The corrected analyses ran
in Python 3.12.9 with PyTorch 2.11.0, transformers 5.5.4, NumPy 1.26.4,
SciPy 1.17.1, pandas 2.3.3 and scikit-learn 1.8.0, as recorded in their
run configuration files; the deployment's runs did not record package
versions. All code is at
\url{https://github.com/Biodyn-AI/maxtoki-mechinterp}.

\subsection*{Declaration of generative AI and AI-assisted technologies}
Generative AI was the object of this study and also a tool in it. The
deployment (specifications, pipeline execution, run summaries and the
in-session audit) was carried out by Claude Opus~4.7 through Claude Code,
except for one automated loop that called Claude Sonnet, as described
above. The controlled studies used Claude Opus~5.5 as subjects
and graders. Opus~5.5 agents also wrote the Study~A answer keys (a
second agent re-derived their key numbers before the keys were frozen),
classified review findings that matched no ledger entry, checked disputed
Study~B findings by running code, and did the post-hoc labelling of false
alarms, as described above. The protocol, the Study~A tasks, the Study~B
packages, the deterministic checker and the supporting files were drafted
by Opus~5.5 agents under the author's direction and checked by separate
agents and by code. The exploratory arm of Study~A used Claude Haiku~4.5 as
executor, reviewer and repairer, with Opus~5.5 graders, as described
above. Claude Opus~5.5 also carried out the independent verification, the
corrected analyses (the key numbers of each were computed a second way
with separate code; other agents also re-derived them, three for the RPE1
attention analysis and one for the cross-model comparison), the statistical
analysis code and the figure code, and
drafted this manuscript under the author's direction. The author reviewed
and edited all content and takes full responsibility for it. AI tools were
not used to create or alter any image of primary data; all figures are
plotted by code from the stored outputs.

\section*{Results}

The results come in two parts. The first part evaluates the framework:
the errors in the deployment's outputs, the blind re-audit (Study~C), the
controlled execution study (Study~A), the held-out benchmark (Study~B)
and what deterministic checks can and cannot catch. The second part
reports the MaxToki-217M case study after correction.

\subsection*{Errors in the deployment's outputs}
The ledger holds 149 confirmed errors in the deployment's code, logs,
outputs, run summaries and audit reports: 13 critical, 57 major and 79
minor (Fig~\ref{fig:ledger}; S1~Table). Of these, 129 were present before
the in-session audit. Two more (L002 and L085) were partly present: an earlier form of each error was in the files before the audit. The audit's repairs then rewrote that
text, and the error is in the rewritten version. The other 18 were introduced
by the audit itself. Sixteen are in files written during its two repair rounds;
one of these (L093) also appears in its sweep report. The other 2 are in
its sweep report (L091) and in the checklist that the agent wrote that
day (L092), both written before the repairs began. Two of the 18 are
critical. One is a transcription-factor ``positive'' for GATA1 that does
not survive a null matched for how often each gene is detected; it was
found by trying four thresholds after the original test failed. The
other is an audit summary that, after repair, called 60 of the 80
entries in its pipeline-by-pattern verdict matrix (75\%) ``clean'', with
no new matrix behind that count.

\begin{figure}[H]
\caption{\textbf{Errors in the deployment's outputs.} (A)~The 149
confirmed errors in the deployment's code, logs, outputs, run summaries
and audit reports, by type and severity (critical: a headline result or
claim cannot stand; major: it stands only with material qualification;
minor). (B)~Who or what found each error first. ``In-session audit''
covers errors first flagged, fully or only partly, in the audit's sweep
report. ``Deployment agent, outside the audit sweep'' covers errors the
executing agent found while running the analyses or during the audit's
two repair rounds. ``Human reader'' is a person from outside the
project, not the author.
Source data: S1~Table; the type used for each error in panel A is listed
in the figure's source data.}
\label{fig:ledger}
\end{figure}

Ten of the 13 critical errors were of four kinds. The other three are
the two critical errors that the audit introduced (above) and the
comparison of the direction-of-change accuracy with 50\% (below).
(i)~\emph{Intervention code that
edited the wrong tensor.} A hook is a short piece of code that changes a
layer's activations while the model runs. Here hooks were used to ablate
a sparse-autoencoder feature (set it to zero) or to boost it. In the
circuit-tracing, exhaustive-mapping and steering scripts, the hook
replaced the output of transformer block~$l$ with the input of that block
plus the feature edit. This deleted the whole block. At the last layer,
the steering hook instead applied the final normalisation twice.
Combined edits also overwrote each other, so ablating three features gave
exactly the effect of ablating the last one. With the feature edit set to
exactly zero, the deployed steering code still gave 98.5\% to 101.9\% of
the effect that the deployment reported for one steered feature at each
of the five layers. It gave 93\% to 101\% of the deployment's mean over
the three steered features of each layer. With a zero edit, the deployed
circuit code gave a set of 71,474 edges, and this set held more than 99\% of the edges of
each of the 120 traced features. So the deployed steering effects and
circuit edges came from the hook bug, not from the features.
(ii)~\emph{Wrong model input.} Every run except the RPE1 attention run
ranked genes by $\log(1+\text{CP10k})$ values ($\log(1+\text{CPM})$,
counts per million, in one external lung dataset) divided by the gene
median, instead of by counts divided by the gene median. On average,
30\% to 51\% of the 200 highest-ranked genes differed from the correct
encoding (dataset means over 35 to 55 cells each), and the position of
nearly every gene changed. (iii)~\emph{Metrics that could not measure
what they claimed.} The CRISPR-interference direction-of-change accuracy
counted every target that went down as correct, whatever the prediction.
The annotation of sparse-autoencoder features applied a false-discovery
correction only to tests that had already passed $p<0.1$. In the
developmental orderings, each cell group's place in the ordering was
looked up from its cell-type label, so a random code for each label
passed the same quality gates. The scGPT cross-model comparison read
another gene's embedding row for every gene. (iv)~\emph{A mislabelled
steering axis.} The steering runs pushed cells along a ``pseudotime''
axis, which should order cells from immature to mature. In fact, the
axis was the first principal component of the expression of 200 randomly
drawn Tabula Sapiens immune cells. Cell type explained 94.5\% of its
variance. The 50 ``earliest'' cells were 49 lymphoid cells (T, B and
plasma cells) and one stem cell. The 50 ``latest'' were all myeloid cells
(neutrophils, monocytes and macrophages). So the axis separated two
blood-cell lineages, not immature from mature cells.

Fig~\ref{fig:ledger}B shows who or what found each error first. The
sweep report of the in-session audit was the first to flag 10 errors (4
major, 6 minor). For 3 of them (1 major, 2 minor) it flagged only the
area, not the error itself. It flagged none of the critical errors. The
deployment agent found 4 more (1 critical, 1 major, 2 minor): one while
running an analysis and three during the audit's two repair rounds. The
critical one was the direction-of-change metric, found in the second
repair round. Independent verification found 71, the blind re-audits of
Study~C 60, and a human reader 4. Of the 13 critical errors, independent
verification first found 8, the blind re-audits 3, the deployment agent 1
and the human reader 1. The one the human reader found was that the
direction-of-change accuracy was compared with 50\%, both before and after
its repair, instead of with the accuracy of always predicting a decrease.
The repair rounds repaired 8 errors that were present before the audit,
3 of them only partly.

The audit's sweep of all eight pipelines took about 8~minutes (3.7~hours
with the two repair rounds). The sweep rated as
clean, or as exemplary, several analyses that later proved to be
artefacts. For example, it called the nulls of the developmental-order
analysis the project's ``gold standard'', although a random code for each
cell-type label passed the same quality gates (S1~Table, R58). It also
rated the degree-preserving null of the attention analysis ``clean ---
exemplary'', although in the MaxToki-217M K562 run 49 of that null's 50
draws were identical to the real reference network (L033).

\subsection*{Study C: blind re-audits find four times more errors than the in-session audit, with or without the checklist}
Nine blind re-audits of the pre-audit snapshot (three per arm) each took
21 to 29~minutes and 84 to 105 tool calls, and returned 36 to 49
findings. Scored against the 68 reference errors, fresh agents fully
detected 43\% of them with the checklist as deployed (95\% bootstrap
interval over errors 32--53\%), 45\% with the generic checklist
(35--55\%) and 41\% with a generic critical review (31--51\%); the sweep
report of the in-session audit had detected 10\% (4--18\%)
(Fig~\ref{fig:studyC}A).
Counting partial detections, the three arms reached 73\%, 71\% and 61\%.
The arms did not differ clearly: the generic checklist detected 4.4
percentage points more than the generic review ($-0.5$ to $+9.3$;
Wilcoxon signed-rank $p=0.10$ over the 20 errors on which they differed),
and the checklist as deployed 2.0 points more ($-5.9$ to $+9.8$;
$p=0.42$). Each fresh arm detected 30 to 35 points more than the in-session
audit.

\begin{figure}[H]
\caption{\textbf{Study C: blind re-audits of the pre-audit snapshot.}
(A)~Share of the 68 reference errors fully detected by each arm (filled
marker; line, 95\% bootstrap interval over errors), detected fully or
partly (open marker), and by each single run (ticks). The in-session
audit is scored from its written report. (B)~Detection of each reference
error (rows, grouped by severity) by the in-session audit (A) and by each
of the nine re-audits (G, generic review; C, generic checklist; D,
checklist as deployed; 1--3, runs).}
\label{fig:studyC}
\end{figure}

The arms differed in what they found. Averaged over its three runs, the
generic review fully detected 81\% of the 12 code bugs, against 69\% for
the generic checklist and 58\% for the checklist as deployed. It also
fully detected five of the six critical errors in every run
(Fig~\ref{fig:studyC}B). The checklist arms did better on two other
types. They fully detected 51\% (generic checklist) and 55\% (as
deployed) of the 17 missing baselines and nulls, against 39\% for the
generic review. They fully detected 22\% and 37\% of the 9 wrong
statistics, against 15\%. Fig~\ref{fig:studyC} does not show error
types; the type of each error is in the figure's source data. Together
the nine audits fully
detected 44 of the 68 errors; 6 (2 major, 4 minor) were never detected,
even partly. One critical error, the scGPT row lookup, was detected only
partly: reviewers saw that the cross-model correlation sat at chance, but
none traced it to the lookup. Only part of the wrong input encoding was
in the reference set: three attention runs fed $\log(1+\text{CP10k})$
values in as if they were counts (S1~Table, L118, major). The full error,
that every run except the RPE1 attention run was encoded on the wrong
scale (L124, critical), was found after the set was frozen. Three audits
were graded as fully detecting L118 and two as detecting it partly. All
five suspected that the Adamson input was log-normalised and noted that
no code checked whether the stored values were raw counts. None found
that the K562 inputs were log-transformed, and none traced the error to
all the runs it affected.

Beyond the reference set, the nine audits raised 132 findings that
matched no ledger entry. The checking agent judged 123 of them real
errors, 6 real but trivial, 2 false alarms and 1 undecidable. After
de-duplication across audits the real errors formed 62 candidate errors,
of which 60 were confirmed against the full project and added to the
ledger. The graders
agreed closely on detection ($\kappa = 0.94$, range 0.86--1.00 across the
nine audits).

\subsection*{Study A: with Claude Opus 5.5 every executor reached the correct verdict with or without the specification, and generic reviews turned three correct verdicts wrong}
All 40 executor runs and all 120 final deliverables were completed and
graded. Every executor run reached the key's verdict, in both arms (20 of
20 with the specification and 20 of 20 without it; risk difference 0,
Newcombe 95\% interval $-16$ to $+16$ percentage points; Fisher
$p = 1$; Mantel--Haenszel $p = 1$) (Fig~\ref{fig:studyA}A). After review
and repair, 40 of 40 deliverables were correct in the checklist arm and 37
of 40 in the generic-review arm (exact McNemar $p = 0.25$ against the
checklist; $p = 0.25$ against no review). None of the four primary
contrasts was significant (Holm-adjusted $p = 1$ for each). The two
graders agreed on the verdict of all 120 deliverables.

\begin{figure}[H]
\caption{\textbf{Study A: controlled execution study with Claude Opus~5.5
as executor, reviewer and repairer.} (A)~Number of the five executor runs
per task whose final deliverable reached the key's verdict, by executor
arm (paper: brief, data and source paper; contract: the same plus the
deployed specification) and by what followed the executor (no review, a
generic critical review then repair, or a generic-checklist review then
repair). (B)~Share of each task's known pitfalls handled correctly (mean
over deliverables; line, 95\% bootstrap interval over executor runs
pooled over tasks).
(C)~Review points judged valid or invalid (or unclear) by a grader with
the key, over the 40 reviews of each style.}
\label{fig:studyA}
\end{figure}

The three verdicts that became wrong were all on the cross-model
alignment task (T4) in the contract arm (where the executor also received
the deployed specification), after a generic review. In each,
the reviewer argued that the executor's decision rule was too lenient, or
that metrics named in the specification (the mean in-sample canonical
correlation, or the Spearman correlation of pairwise distances) gave a
weaker result than the executor's chosen metric. The repairing agent
accepted these points, rewrote its decision rule, and changed the verdict
from ``aligned with scGPT but clearly less than with Geneformer'' to
``inconclusive''. The key does not accept ``inconclusive'': every estimate
is far above chance, and the Geneformer-minus-scGPT difference has an
interval that excludes zero on every panel. Canonical correlation finds
pairs of directions, one in each embedding table, along which the genes'
positions correlate as strongly as possible. When it is fitted and scored
on the same genes (in-sample), as the reviewers used it, it is high even
for shuffled genes.
In the setting these executors used (20 principal components, mean of all
20 canonical correlations), their own permutation tests put chance at
about 0.20. The observed values were 0.40 for scGPT and 0.44 for
Geneformer. This chance level depends on the number of components as well
as on the number of genes, and the specification does not state it.
Graders with the key judged 17 of the 22 points in these three reviews
invalid. The checklist reviews of the same three executor outputs
did not change the verdict.

The secondary outcomes moved little (Fig~\ref{fig:studyA}B). Executors
handled most of each task's pitfalls without help (94.5\% on average
without review). Review and repair raised this slightly: by 2.8
percentage points with the checklist (95\% bootstrap interval 0.9 to 4.7)
and 2.0 points with a generic review (0.3 to 3.7); the checklist and the
generic review did not differ ($+0.8$; $-0.9$ to $+2.5$). Without review,
executors with the specification handled 2.3 points fewer pitfalls than
those without it ($-5.4$ to $+0.8$) and reported 9.2 points more key
numbers within tolerance ($-0.6$ to $+21.3$). These intervals resample
executor runs pooled over tasks. The pre-registered bootstrap resamples
runs within each task. It gives the same conclusions for the effects of
review on pitfalls handled, in percentage points: 1.1 to 4.4 for the
checklist against no review, 0.5 to 3.4 for a generic review against no
review, and $-0.9$ to $+2.4$ for the checklist against a generic review.
It also gives the same conclusions for the effects of review on pitfalls
handled in the Haiku arm below. But for the two differences between the
executor arms, its intervals just exclude zero: $-4.4$ to $-0.2$ points
for pitfalls handled and $+0.6$ to $+18.9$ points for key numbers within
tolerance. Both
differences are small, and whether their intervals exclude zero depends on
how the runs are resampled. Materially false statements were rare (0.13
per deliverable) and did not differ between conditions.

Most review points were judged invalid by a grader with the key, meaning
the point was wrong, already handled in the deliverable, or did not
matter for its results: 61\% of the 281 generic-review points and 70\% of
the 297 checklist points (208 invalid and 1 unclear)
(Fig~\ref{fig:studyA}C). Only in the three T4 cases above did acting on
such points turn a correct verdict into a wrong one. Three other repairs,
all on T1 in the paper arm (two after a generic review and one after a
checklist review), changed ``not supported'' to ``inconclusive''. The T1
key accepts that answer under stated conditions, so all three stayed
correct.

An executor run took a median of 27~minutes and 27 tool calls; a review
took a median of about 14~minutes and 16 (checklist) or 18.5 (generic)
tool calls, and a repair 21 to 25 tool calls.
All 120 returned analysis scripts ran on the task data and printed their
results. Six of them were stopped at least once because the shared
machine was busy: five ran past the 20-minute limit, and two were stopped
without output (one script had both). Run again with fewer scripts at a
time, or on their own, all six finished in under 10~minutes. Of the
11,045 estimates that the deliverables reported, 10,363 (94\%) could be matched by name to a number
the script printed, and all 10,363 agreed within tolerance. Of the 682
unmatched estimates, 681 were in T3, whose scripts printed their results
under other names. For these, the only check left was to look for each
value among all the numbers a script printed, within the same tolerance,
and that check says little. In the 21 T3 scripts concerned, 98--100\% of
the deliverable's own estimates were found this way, but so were
61--97\% (median 89\%) of the same number of estimates drawn at random
from deliverables of other tasks (mean of 50 draws per script).

In the exploratory arm with Claude Haiku~4.5 as executor, reviewer and
repairer (S2~Appendix), executors without review reached the key's verdict
in 15 of 20 runs with the specification and in 10 of 20 without it
(difference $+25$ points, Newcombe 95\% interval $-5$ to $+49$; Fisher
$p = 0.19$); most of the difference came from T1 (3 of 5 against 0 of 5).
Haiku executors handled far fewer pitfalls than Opus executors (30--36\%
against 93--96\% without review). A checklist review followed by repair
raised the share of pitfalls handled by 9.2 percentage points (95\%
bootstrap interval over executor runs pooled over tasks, 4.9 to 13.6;
pre-registered interval resampling runs within each task, 5.0 to 13.3).
It raised it by 6.4 points more than a generic review followed by repair
(over the 28 runs with both reviews valid: 1.3 to 11.7 pooled over tasks,
2.1 to 10.4 within tasks). But it did not
raise the share of correct verdicts. Of the 33 executor runs with a valid
checklist review, repairs turned 3 of 12 wrong verdicts right and 3 of 21
right verdicts wrong. Of the 35 executor runs with a valid generic review,
repairs turned 2 of 13 wrong verdicts right and none of 22 right verdicts
wrong. Twelve reviews (7 checklist, 5 generic) returned no valid output,
so 12 of the 120 deliverables are missing.

\subsection*{Study B: on held-out analyses both review styles find every documented error, and both over-call errors in correct analyses}
In 132 reviews of the 22 held-out items, every review of a flawed item
found its documented error, in both arms. This held for all 11 flawed
items and for 33 of 33 reviews per arm (Wilson 95\% interval over items
74--100\%; items are the unit because the three reviews of an item are
not independent). It held whether the error fell inside the
checklist's patterns (6 items) or outside them (5 items: a circular
correlation that fails on evenly spread angles, span-filled residue
labels, an arbitrary sign flip of singular vectors, raw counts fed to an
encoder that expects binned values, and a row-index join that matched the
wrong cells). Both graders judged all 66 reviews of flawed items to have
found the documented error, so they agreed on all 66. Cohen's $\kappa$
is not defined when every judgment is the same, so we report this raw
agreement instead.

On the 11 correct (clean) items, most reviews still claimed at least
one error in a main result that the item's key does not count as an
error (Fig~\ref{fig:studyB}A). This happened in 79\% of checklist
reviews (26 of 33; 95\% bootstrap interval over items 58--94\%) and in
82\% of generic reviews (27 of 33; 61--100\%). The paired difference,
checklist minus generic, was $-3$ percentage points (95\% bootstrap
interval over items $-21$ to $+15$; Wilcoxon $p=0.88$). Checklist
reviews made 45 such claims in all (1.4 per review of a clean item), and
generic reviews made 43 (1.3 per review). Checklist reviews were longer
(9.2 against 7.8 findings per review); both arms found about two real
problems per review beyond the documented error (2.1 in each).

\begin{figure}[H]
\caption{\textbf{Study B: held-out analyses from other projects.} Every
review detected the documented error in every flawed item, in both arms
(33 of 33 each; not plotted). (A)~For each of the 11 correct (clean)
items, the number of its three reviews per arm that claimed at least one
error in a main result that the item's key does not count as an error.
(B)~Post-hoc classification of the 82 such claims that the graders
marked as false alarms, on clean and flawed items, by arm. The 23
further claims that the verifier ruled false alarms after a grader
dispute are not included.}
\label{fig:studyB}
\end{figure}

Across clean and flawed items, the reviews made 105 such claims (88 on
clean items, 17 on flawed items). The graders marked 82 of them as false
alarms (73 on clean items, 9 on flawed items). The other 23 (15 on clean
items, 8 on flawed items) were findings that a grader marked as disputed
and the verifier agent then ruled false alarms. After the results were
known, we classified the 82 claims that the graders marked. This step was
not pre-registered. One agent labelled each claim, with no second rater.
Most of the 82 were accurate observations whose severity was overstated
(Fig~\ref{fig:studyB}B): 59 were correct but did not change any stated
conclusion (for example, that an input-label overlap weakened, but did
not overturn, an interpretation), 17 were correct and already stated in
the analysis's own summary or README (for example, that a replication in
a second cell line rested mostly on one transcription factor), 3 were
factually wrong and 3 were real problems that the key had missed. The
pattern was the same in both arms. The 23 claims ruled by the verifier
were not classified.

\subsection*{What deterministic checks catch}
Run on the eight deployed specifications, the specification linter
rejected seven: all except the ageing-probe specification lacked the
code-references section (so their later sections were also numbered out
of template order), and only 5 of 230 code pointers used the required
\texttt{repos/<name>/path:LINE} form. Six of the seven pinned reference
repositories (the manifold specification names none) had been copied
without their version-control history, so their pins could be matched
only as text. The run-manifest check found 10 parameters that the runs
set differently from their specifications (for example, 50 instead of
200 draws of the degree-preserving null, and 50 instead of 200 cells per
ablation condition); a written reason existed for one of them. One of
the 10, the number of control cells, differed only in the attention run
on MaxToki-1B. The check could compare only values that a run had
recorded: 66 of the 258 parameter entries in the specifications. The
other 192 could not be checked, so further differences may have been
missed. The check also found two manifold quality-gate verdicts recorded
as passed although one of the gates they required had never been
computed. The run-manifest check also found that all 75 deployment
scripts that use random numbers set a seed in the same script. No run
wrote a run manifest, and 3 of the 8 runs (the ageing probe, circuit
tracing and exhaustive mapping) recorded no device. No run recorded its
package versions or which MaxToki revision it loaded (see Materials and
methods). The package list with pinned versions was saved during the
in-session audit, after most runs. Some manifold and topology steps used
a second Python environment whose packages were not recorded (S1~Table,
L082, L088 and L089).

The number tracer found 1,132 of 1,186 distinct numbers in the run
summaries (95.4\%) in their run's outputs. Numbers of the same format
drawn at random were also found 78\% of the time, so tracing exceeded
chance by 18 percentage points (95\% interval over runs 8--27); the
excess was smallest in runs with many numeric outputs. A traced number
can still be described wrongly. For example, the manifold intervals were
traced correctly to their outputs, and the run summary said how they were
made: 200 repeats of 80\% subsampling of cell groups, without
replacement. But the summary called them bootstrap intervals. It presented them as the
fix the audit had asked for, which was resampling at the unit of inference
(here the donor). It then used them to call the gate passes robust to
resampling of cell groups. Cell groups
from the same donor are not independent, so these intervals are too
narrow. Resampling donors widens the interval for one gate from
0.270--0.480 to 0.142--0.540, so its lower end falls below the 0.20 pass
mark (S1~Table). The checks
can establish that a document has the required structure, that a
parameter differs from its default and that a number appears in an
output. They cannot establish that a null randomises, that an
intervention edits the right tensor, that the model input is encoded as
the model expects, that a lookup reads the right rows or that a number
means what the text says. Twelve of the 13 critical errors in the ledger
were of this second kind. The thirteenth was an audit summary that,
after repair, called 60 of the 80 entries in its pipeline-by-pattern
verdict matrix (75\%) ``clean'', with no new matrix behind that count
(S1~Table, L090). A number tracer looks for this kind of problem. But we
ran the tracer only on the run summaries, not on the audit reports, so we
did not test whether it would have flagged it.

\subsection*{The MaxToki-217M case study}
This part reports what the corrected analyses show about MaxToki-217M.
Every result concerns MaxToki-217M given one cell at a time, with the
input encoding and the intervention code checked as described in
Materials and methods. Table~\ref{tab:scope} lists, for each analysis,
the cells it used, the reference it was judged against, and what its
conclusion depends on. Four deployed analyses used incorrectly encoded
inputs and were not re-run, so we do not report them: the topology
screen of gene--gene representations, the ageing probe, the
feature-patching specificity ratio, and the edge census from 1,000
layer-5 features in 20 cells, which also used the intervention code that
deleted a transformer block. The MaxToki-1B attention run and the
Adamson attention run are not reported for the same reason. The
gene-set annotation of the sparse-autoencoder features was redone only
to choose the source features for circuit tracing, and its rates are not
reported. S3~Appendix lists smaller follow-up analyses that were not
re-run either, such as value-weighted attention, head ablation and the
later uses of the circuit edge list.

{\footnotesize
\setlength{\tabcolsep}{3pt}
\renewcommand{\arraystretch}{1.15}
\begin{longtable}{@{}p{0.15\linewidth}p{0.20\linewidth}p{0.18\linewidth}p{0.20\linewidth}p{0.24\linewidth}@{}}
\caption{\textbf{Scope of the MaxToki-217M results.}}
\label{tab:scope}\\
\toprule
\textbf{Analysis} & \textbf{Cells} & \textbf{Judged against} &
\textbf{Conclusion depends on} & \textbf{Result} \\
\midrule
\endfirsthead
\multicolumn{5}{@{}l}{\footnotesize\emph{Table \thetable\ continued from previous page.}}\\
\toprule
\textbf{Analysis} & \textbf{Cells} & \textbf{Judged against} &
\textbf{Conclusion depends on} & \textbf{Result} \\
\midrule
\endhead
\midrule
\multicolumn{5}{r@{}}{\footnotesize\emph{continued on next page}}\\
\endfoot
\bottomrule
\endlastfoot
Attention, RPE1 & 2,000 RPE1 control cells; layer~8; 1,500 genes &
RPE1 knockdown responses (325 knockdowns); TRRUST (16 factors) &
Causal attention mask; few TRRUST factors per cell line &
Below the target-variance baseline; symmetric score beats the network
null but not co-expression \\
Attention, K562 & 2,000 K562 control cells; layer~8; 1,500 genes & K562 knockdown responses (194 knockdowns); TRRUST (16 factors) & Causal attention mask; one factor (MYC) carries the regulator-to-target signal; which genes enter the test & At chance on knockdowns; ties co-expression on TRRUST; no score beats the network null after correction \\
Factor-specific SAE features & 9,200 K562 cells (87 knocked-down
factors); layer-5 SAE & DoRothEA ChIP-seq targets (20 factors); TRRUST &
SAE dictionary; top-20-gene summary of a feature; target-set coverage &
No factor passes after Benjamini--Hochberg correction under four nulls;
16 factors pass a fifth null that does not match the number of
responding features. GATA1: a feature
with 8 of 20 top genes among its targets, as the deployment reported, is
ruled out; one with 4 of 20 is not \\
Circuits and knockdown direction & 200 K562 control cells; 120 source
features; 248 silenced genes & K562 CRISPR-interference fold-changes &
SAE features at four layers; edge thresholds; noisy fold-changes &
No skill beyond constant rules; below a no-model rule by 8.3 points \\
Feature triplets & 20 K562 control cells; 512 triplets & Additive
model; planted interaction & Small edits (about 2\% of the residual
norm) & Additive; no three-way interaction \\
Steering & 50 Tabula Sapiens bone-marrow HSCs from two held-out donors; layers 0, 3, 6, 9 and 11 & Erythroid maturity score checked on held-out donors; 20 random features per layer & Sparse autoencoders trained on K562 cells; a score built from the model's output; one lineage; 47 of 50 cells from one donor & 14 of 15 candidates no better than random; one feature that marks erythroid cells moves the score 3.5\% of the span \\
Developmental order & Tabula Sapiens immune cells: 290 cell groups (11 donors) to fit, 600 and 160 groups from other donors to test & Curated 34-stage blood tree; structured null that keeps cell-type classes; cell-type label code, gene expression and input-token features & Deployed read-out head and pass marks; only the T lineage has graded stages; one head seed for most numbers & No order beyond cell type; in cross-lineage order, a small lead over one bag of input tokens on the zero-shot panel (a tie on the external panel over head seeds) and a tie with a second bag built from the deployed gene order \\
Spectral geometry & Three disjoint samples of 2,000 Tabula Sapiens immune cells; 1,500 genes; all layers & Three randomly initialised copies of the model; split halves of each sample & Per-gene averages over cells; one atlas and gene panel; default random initialisation & Effective rank falls with depth with random weights too; gap of cross-sample CKA from 1 is mostly sampling noise \\
Cross-model alignment & None (input embedding tables); three panels of
350--382 genes & scGPT and Geneformer~V2-316M tables; shuffled genes &
Input tables only, not internal layers & Aligned with both; more with
Geneformer \\
\end{longtable}
\begin{flushleft}
Every analysis gave
the model one cell per sequence, ranked by counts divided by each gene's
median (rank-value encoding). HSC, hematopoietic stem cell; SAE, sparse
autoencoder; CKA, centred kernel alignment, a similarity of two sets of
gene vectors (1 means the same geometry). The span is the mean maturity
score of erythrocytes minus that of HSCs. A bag of input tokens describes
each cell group by the genes in its cells' model input and uses no model.
\end{flushleft}
}


\subsubsection*{Attention scores do not beat gene-level baselines}
MaxToki is a causal decoder: genes enter in order of falling normalised
expression, and a gene can attend only to genes placed before it. The
attention from a regulator $P$ to a target $T$ is therefore zero in every
cell in which $T$ comes after $P$, and the average over cells partly
measures how often $T$ is ranked above $P$ (the order share), which
needs no model. We read five attention scores at layer~8: regulator to target, target to
regulator, the mean and the larger of the two directions, and regulator
to target averaged only over the cells in which the target comes before
the regulator. A gene counted as responding to a knockdown when its mean
$\log(1+x)$ value in the knocked-down cells differed from that in 5,000
control cells by at least 0.5 (natural log), with BH $q < 0.05$ over the
1,500 genes tested (Welch $t$-test). Here $x$ is the count scaled so that
each cell sums to 10,000 over those 1,500 genes. A knockdown was kept
when it had at least 3 responding and 3 non-responding genes. The
deployed RPE1 pipeline used this rule. We kept it for K562 and chose it
before seeing the K562 results. On the RPE1 screen, which was encoded
correctly from raw counts, the order share alone predicted knockdown responses better than
the regulator-to-target score (mean AUROC over 325 knockdowns 0.648
against 0.603), and the variance of the target gene predicted them far
better (0.766, 95\% bootstrap interval over knockdowns 0.757--0.776)
(Fig~\ref{fig:caseReg}A). AUROC is the probability that a randomly chosen
responding gene is scored above a randomly chosen non-responding one;
0.5 is chance.

\begin{figure}[H]
\caption{\textbf{MaxToki-217M: regulatory probes after correction.}
(A)~Prediction of which genes respond to a CRISPR-interference knockdown,
in RPE1 cells (325 knockdowns) and K562 cells (194 knockdowns): mean
AUROC over knockdowns (line, 95\% bootstrap interval over knockdowns) for
layer-8 attention read from regulator to target and in both directions,
and for three scores that use no model: the share of cells in which the
target is ranked above the regulator, co-expression, and the variance of
the target gene. Dashed line, chance. (B)~Recovery of TRRUST
regulator--target pairs (RPE1: 16 factors, 138 pairs; K562: 16 factors,
107 pairs): pooled AUROC with 95\% bootstrap interval over factors; the
target-gene model uses each target's mean, variance and dropout rate.
(C)~Power of the count-matched test for GATA1 in K562 cells, counting
every overlapping target (the deployed analysis stopped at 5): share of
200 simulations in which GATA1 ChIP-seq targets planted among the top 20
genes of one responding feature were detected at $p<0.05$ (band, Wilson
95\% interval). (D)~Prediction of the direction of change of 796,304
knockdown--target records in K562 cells (248 silenced genes) by the
traced circuit (line, 95\% bootstrap interval over silenced genes), by
always predicting ``down'', and by predicting each target's usual
direction from other silenced genes.}
\label{fig:caseReg}
\end{figure}

On TRRUST regulator--target pairs in RPE1 (16 transcription factors, 138
pairs), a symmetric attention score (the mean of both directions) reached
an AUROC of 0.638 and beat a degree-preserving network null ($z = +3.2$;
empirical $p = 0.003$ from 1,000 rewired networks)
(Fig~\ref{fig:caseReg}B). After Benjamini--Hochberg (BH) adjustment the
$q$ value was 0.018. The adjusted family had 12 tests: six scores (the
five attention scores and the order share) in each of the two cell
lines. Co-expression computed from the same cells beat the null about as
strongly (AUROC 0.596, $z = +3.0$). Attention scored 0.042 higher, but
the 95\% interval of this gap ($-0.040$ to $+0.127$) includes zero. A
three-number model of each target gene knows nothing about the
regulator. It uses the gene's mean, its variance and its dropout rate
(the share of control cells in which the gene has zero counts). It scored
0.749 (0.691--0.801). The symmetric attention score was 0.11 below this
gene model (0.00 to 0.20). We also removed the mean, variance and
dropout rate of both genes from the symmetric score with the planned
linear model. What remained had an AUROC of 0.538 (0.456--0.614), which
does not differ from 0.5. It still beat the rewired networks
($z = +2.75$, $p = 0.003$, not adjusted). This is because the rewired
networks were scored after the same removal, and their mean AUROC fell
too, from about 0.58 to about 0.48. Co-expression treated the same way
also beat them ($z = +2.94$), and the symmetric score still beat them
after co-expression was removed
(0.611, $z = +2.53$). So in RPE1, symmetric attention carries a
pair-level signal about as large as that of co-expression, not a larger
one (S3~Appendix).

In K562 cells (2,000 control cells of the same screen, 1,500 genes,
layer~8), attention predicted knockdown responses at chance under this
rule. Over 194 knockdowns the mean AUROC was 0.505 (95\%
bootstrap interval over knockdowns 0.491--0.520) from regulator to target
and 0.502 (0.488--0.517) for the symmetric score; all five attention
scores lay between 0.499 and 0.509 (Fig~\ref{fig:caseReg}A). Target
variance scored 0.690 (0.674--0.708) and co-expression 0.556
(0.540--0.573). A second label rule ran the same test on the stored
$\log(1+\text{CP10k})$ values. These are scaled over all 6,546 genes in
the file, not over the 1,500 genes. Under this rule only 67 knockdowns
had at least 3 responding and 3 non-responding genes. With it
the regulator-to-target score was 0.550 (0.517--0.580). That is above
chance, but below the order share (0.611, 0.568--0.651), which needs no
model, and below target variance (0.695, 0.659--0.730) (S3~Appendix).
On the 107 TRRUST pairs of 16 factors in this gene set, the symmetric
score had a pooled AUROC of 0.571 (0.521--0.617) and $z = +1.76$ against
the rewired networks (empirical $p = 0.041$; BH $q = 0.069$ over the 12
tests), so the RPE1 result did not replicate in K562
(Fig~\ref{fig:caseReg}B). The regulator-to-target score had a pooled
AUROC of 0.548 (0.499--0.594) and $z = +1.72$ ($q = 0.069$), and it
rested on one factor. MYC holds 27 of the 107 pairs, and without MYC
this $z$ fell to $+0.12$ (and the symmetric score's $z$ to $+1.11$).
Co-expression from the same cells did as well as attention (0.588,
$z = +1.55$), and the target-gene model did better (0.723,
0.644--0.814). After the same gene-level features were removed, no K562
score beat the rewired networks (lowest BH $q = 0.18$ over the six
scores with the planned linear model, and 0.17 with a boosted-tree
model). Which genes enter the test matters. Our gene set is the 1,500
genes with the highest variance of log-normalised expression. The
deployed analysis used the same rule on values that had been logged a
second time, and only 627 of its 1,500 genes are in our set. On the
deployed gene set (8 factors, 45 pairs), the symmetric score did beat
the null ($z = +2.38$, $p = 0.010$). This check was not planned and is
not adjusted. Single heads, other layers as a primary endpoint and
references other than TRRUST were not tested (S3~Appendix).

\subsubsection*{No sparse-autoencoder feature is specific to a knocked-down transcription factor}
A sparse autoencoder (SAE) rewrites a hidden state as the sum of a few
learned directions, called features, out of a much larger set. We
trained twelve SAEs on 500 K562 control cells, one for each hidden state
the model exposes. Layer~0 is the token embedding. Layers 1--10 are the
outputs of the first ten of the 11 transformer blocks; this hidden state,
passed from one block to the next, is called the residual stream.
Layer~11 is the output of the last block after the final normalisation.
Each SAE has 4,928 features and keeps only the 32 most active features
for each token (a TopK SAE). The SAEs explained 70\% of the held-out
variance at layer~0 and 84--93\% at layers 1--11. For 87 transcription
factors knocked down in K562 (34 to 100 cells each; 20 with ChIP-seq
target sets in DoRothEA~\cite{garcia-alonso2019dorothea}), we called a
layer-5 feature responsive when its activation differed from control
cells (Mann--Whitney test, BH $q < 0.05$, and a minimum effect size),
and asked whether the 20 genes on which a responsive feature is most
active overlap the factor's targets more than expected. No factor passed
after BH correction across factors, at any of seven effect cut-offs and
under any of four nulls. In the count-matched null, each top gene is
swapped for a random gene seen in a similar number of cells. The second
null also matches gene length. The third uses random groups of control
cells that keep as many responding features as the knockdown. The fourth
uses the responding features of other knockdowns. The smallest $q$ was
0.18 (GATA1 at the 0.5 cut-off, under the other-knockdown null, where
the smallest possible $p$ is 0.059). We also ran a fifth null. It uses
random groups of control cells of the knockdown's size and runs the
feature selection again on each group. Unlike the third null, it does
not keep as many responding features as the knockdown. Every factor
with at least 2 targets among a responding feature's top 20 genes passed
that null (49 tests, 16 factors). This happens because random control
groups rarely select any feature, and never at cut-offs of 0.25 or
above. So that null shows only that the knockdown changed something, not
that the features know the factor's targets (S3~Appendix).

GATA1 is the factor the deployment singled out. At the cut-off the
deployment used (a difference of more than 0.5 between knockdown and
control cells in a feature's mean activation), seven features responded
to the knockdown. The best had 4 of its 20 top genes among the
221 GATA1 ChIP-seq targets in the gene set, where 3.1 are expected by
chance ($p = 0.32$ under the count-matched null, $0.43$ when gene length
is also matched, $0.19$ against random control groups with the same
number of features); GATA1 ranked 82nd of 291 factors for these
features. The test detects strong but not weak specificity. We planted
GATA1 targets among the top 20 genes of one responding feature, with 200
simulations for each number of planted targets. With 6 or more planted
targets, the count-matched test gave $p < 0.05$ in all 200 simulations.
With 5 it did so in 80\% of simulations, and with 4 in 29\%
(Fig~\ref{fig:caseReg}C). These power values count every overlapping
target. The deployed statistic stopped counting at 5. For GATA1 that
version gives the same $p$ here (0.32), but its smallest possible $p$ is
0.099, so it could not detect any planted signal at this cut-off
(S3~Appendix). The data thus exclude a GATA1 feature as specific as the
one the deployment reported (8 of 20), but not a weaker one.

\subsubsection*{Traced circuits do not predict the direction of knockdown responses}
Ablating 120 source features (30 at each of layers 0, 3, 6 and 9) in 200
K562 control cells and reading the change in every later feature gave
318,387 edges between features. A source--target pair counted as an
edge when the target's change had a Cohen's $d$ (its mean change over the
200 cells divided by the standard deviation of that change) above 0.5 in
absolute value, and went the same way in more than 70\% of cells. In
57.9\% of edges, removing the source lowered the target; in the rest it
raised it. We turned these edges into predictions about genes. Each
feature has 10 top genes, the genes on which it is most active. Each edge
pairs every top gene of its source feature with every top gene of its
target feature. We kept a gene pair if at least 2 edges linked it, or one
edge with $|d| > 2$. We predicted that silencing the first gene lowers the
second if more than half of the edges linking them were ones in which
removing the source lowered the target, and that it raises the second
otherwise. We checked these predictions in K562 cells in which the first
gene was silenced by CRISPR interference. The observed direction is the
sign of the second gene's mean $\log(1+\text{CP10k})$ value in the
silenced cells minus its mean in 3,000 control cells. This gave 796,304
records from 248 silenced genes; each record is one pair of a silenced
gene and a measured target gene. The prediction was right
50.3\% of the time (95\% interval over silenced genes 49.9--50.7\%),
below the 52.5\% obtained by always predicting ``down''
(Fig~\ref{fig:caseReg}D). Balanced accuracy, which is 50\% for any
constant rule, was 50.2\% ($+0.19$ points, $-0.18$ to $+0.57$), and the
Matthews correlation was $+0.004$ ($-0.004$ to $+0.011$). A rule that
uses no model, predicting for each target the direction in which it
usually moves when other genes are silenced (fitted on other silenced
genes), was right 58.5\% of the time, 8.3 points more than the circuit
(7.3--9.2). An exploratory split by the size of the observed change found
a balanced accuracy of 51.8\% for changes of 0.1 to 0.25 log-fold, but the
no-model rule reached 68.4\% on the same records.

\subsubsection*{Removing three features at once had an additive effect in the 512 triplets tested}
In 512 triplets of features (eight each at layers 0, 5 and 9) ablated
together and separately in 20 K562 control cells, no triplet had a
three-way interaction at any of the 4,928 last-layer features after BH
correction (2,523,136 tests). The 24 features were drawn at random before
any ablation: at each layer, one from each eighth of activation
frequency, among the features active in at least 1\% of tokens (1,750,
460 and 618 candidates at layers 0, 5 and 9). No annotation was used to
choose them. After subtracting the part expected from cell-to-cell
noise, the share of the joint effect's energy (its squared size) that
was not explained by the three single effects and their pairwise
interactions was 0.0003\% (median over the 512 triplets). Its 95\% basic
bootstrap interval, resampling the 20 cells (1,000 resamples), was
$-0.09\%$ to $+0.06\%$. The interval can go below zero because of the
noise subtraction. The test detected a planted interaction of 5\% of each
cell's joint effect (about 0.2\% of its energy) in 64 of 64 triplets.
Among the 192 feature pairs in these triplets, removing both features
changed the last-layer features by less energy than the two single
changes added together in 5 pairs, and by more in 12 (95\% interval
excluding equality). About 5 per side would be expected by chance, so
these are weak signs. The two clearest pairs both link a layer-5 feature to a layer-9 feature,
with many targets significant after BH correction. In both, removing the
two features together changed the last-layer features by 3--4\% less
energy than the two single removals added together. This fits the
layer-5 feature driving the layer-9 one. Removing all three features together moved the
last-layer residual stream by about 2\% of its norm (median 2.2\%, range
1.6--4.1\%, over 16 cell--triplet pairs), so this describes the model's
response to small edits. The change from single removals was not
recorded. The result covers 8 of the 4,928 features at each of three of
the twelve layers, one read-out layer (layer~11), one cell type and one
SAE training seed. It does not show that three-way interactions are
absent elsewhere in the model.

\subsubsection*{Features that track erythroid maturity move the model's output no further than random features, apart from one feature that marks erythroid cells}
Steering here means scaling one sparse-autoencoder feature up or down
inside the model while it reads a cell, and measuring how far the model's
output moves along a maturity score. The deployment's steering effects
came from the hook that deleted a block, not from the features (see
above). We fixed the cells, score, features, null and tests before any
forward pass. The maturity axis is the erythroid branch of blood
development in Tabula Sapiens bone marrow: hematopoietic stem cells
(HSCs), common myeloid progenitors, erythroid progenitors and
erythrocytes, as labelled in the cell-type annotation. For each cell,
the model's output here is its next-gene scores (logits: one score for
each of the 20,275 tokens in the vocabulary), averaged over all
positions of the cell's sequence. The maturity score is the cosine
similarity of this output with the mean output of erythrocytes, minus
its cosine similarity with the mean output of HSCs. The two means came
from 90 erythrocytes and 76 HSCs of three donors. In cells of two other
donors the score separated erythroid progenitors from HSCs with an AUROC
of 0.991, and its Spearman correlation with stage was 0.84 (95\% bootstrap interval over
cells 0.78--0.88). We give each effect as a share of the span, the mean
score of erythrocytes minus that of HSCs (0.198). The score describes the
model's output; no change in a cell's real state was measured.

We steered 50 HSCs from the two held-out donors (47 of them from one
donor). At each of layers 0, 3, 6, 9 and 11, the three features whose
activity tracked stage most closely were the candidates. Each was
multiplied by 0, 2 or 5 and compared with 20 random features of the same
layer with similar activity. Fourteen of the 15 candidates moved cells no
further in their predicted direction than the random features
(directional $p$ 0.19--1.00; 0.0004\% to 1.5\% of the span at five-fold
strength). The exception, feature 2903 at layer~9, beat all 20 random
features at every strength ($p = 1/21$, the smallest possible; about 0.7
of 15 candidates would reach this by chance) and moved HSCs by 3.5\% of
the span at five-fold strength (95\% bootstrap interval over cells
1.9--5.2\%). It is active at 99\% of gene positions in erythroid
progenitors and at 11\% in HSCs, so it marks the erythroid state. Per
unit of edit it moved cells 37\% as far as a size-matched edit along the
mean difference between erythrocyte and HSC hidden states, and the genes
it raised most (AHSP, ALAS2, HBM, GYPA, HBG1, SLC4A1) are among the ten
that this mean-difference edit raised most. These three checks (where
the feature is active, the comparison per unit of edit and the
most-raised genes) were not planned. Two tests fixed in advance looked at
each layer's three candidates as a group. Neither found them stronger
than random features. First, at no layer did the three candidates push
cells further in their predicted direction than the 1,771 possible sets
of three of that layer's 23 tested features (directional $p$ 0.09--0.64
over the five layers and three strengths). Second, summed over the five
layers, the candidates ranked no higher among the 23 tested features of
their layer than chance (one-sided permutation $p$ 0.24--0.36 over the
three strengths; 100,000 permutations within layers). A third test, also
pooled over layers, found a weak trend. Within each layer, features whose activity rose
more with stage pushed cells a little further toward erythrocytes when
amplified, and a little further back when removed. The mean of the five
within-layer Spearman correlations, each over 23 features, was
0.18--0.19 in size and had the predicted sign (one-sided permutation $p$
0.022--0.031; two-sided 0.045--0.063). Twenty of each layer's 23
features are random features, so this trend does not single out the
candidates. These sparse autoencoders were trained on K562 cells and explained only
34--59\% of the hidden-state variance of the steered HSCs. So single
features of these autoencoders are not maturity switches. This does not
show that the model lacks a maturity direction: adding the full mean
difference moved all cells toward erythrocytes. Other lineages, feature
combinations, autoencoders trained on bone-marrow cells and MaxToki-1B
were not tested (S3~Appendix).

\subsubsection*{Hidden states show no blood developmental order beyond cell-type identity}
The deployed analysis fitted a small read-out head that maps the hidden
states of cell groups (anchors: cells that share donor, tissue, cell type
and stage label) onto a curated tree of 34 blood developmental stages,
and reported a developmental manifold that transfers to new donors. It
judged the head with four pass marks. The first is trustworthiness, at
least 0.80. It checks that anchors that are near neighbours in the
head's output were also close in the head's input, and it does not use
the stage tree. The other three are Spearman correlations between
output distances and stage-tree distances, each at least 0.20. They use
anchors the head was not fitted on, held out at random, by donor or by
lineage branch. The last of these is the branch mark. MaxToki passed
all four on the fitting panel (290 anchors from 11
donors) and on the external test panel (600 anchors from 13 other
donors). The deployment also built a second test panel, which it called
zero-shot: 160 other anchors whose 12 donors are all external-panel
donors, so the two test panels are not independent. On this panel
MaxToki failed the branch mark ($0.147 \pm 0.048$, mean $\pm$ SD over 20
head seeds; 3 of 20 seeds pass). This failure depends on the anchor
set. Some anchors hold only Smart-seq2 cells (Smart-seq2 is a
plate-based sequencing method). When these anchors are dropped (43
fitting-panel, 13 external and 8 zero-shot anchors) and the head is
fitted again, the zero-shot branch mark is $0.34 \pm 0.01$ (mean $\pm$
SD over 10 head seeds; 10 of 10 seeds pass). A random 64-number code for
each cell-type label, which knows nothing about development, passed the three
correlation marks on all three panels. The branch mark can be high just
because anchors with the same label sit together, and only the T lineage
has graded stages. Against a structured null that shuffles stages
between cell-type classes within each branch, MaxToki was not above
the null on the branch mark on any of the three panels. This held both
over all anchor pairs and over only pairs of anchors with different
cell types. MaxToki was also not above the null on cross-lineage
(global) order, the Spearman correlation over all anchor pairs of a test
panel. Across these numbers the one-sided $p$ ran from 0.15 to 0.76 (20
to 2,000 null draws). On pairs of anchors with different cell types, the
branch mark was 0.067, 0.006 and 0.016 on the fitting, external and
zero-shot panels. That is close to zero and well below the
0.20 pass mark. Held-out order within the T lineage was 0.038.

MaxToki was ahead of simple features only in global order
(Fig~\ref{fig:caseGeom}D): 0.865 on the external panel and 0.867 on the
zero-shot panel. That is above the cell-type label code by 0.064 (95\%
donor bootstrap interval 0.050--0.091) and 0.091 (0.027--0.163), and
above expression of highly variable genes by 0.102 (0.073--0.140) and
0.110 (0.038--0.151). A bag of input tokens uses no model. It describes
each cell group by the genes in its cells' model input, and genes nearer
the front of the input get more weight. MaxToki's lead over this bag was
small: 0.016 (0.002--0.046) on the external panel and 0.048
(0.009--0.149) on the zero-shot panel. Over 10 head seeds the lead was
$0.004 \pm 0.005$ on the external panel, a tie, and $0.053 \pm 0.005$ on
the zero-shot panel (above 0 in 10 of 10 seeds). A second bag, built the
same way from the gene order that the deployed analysis fed the model,
also uses no model, and it matched MaxToki on both panels: MaxToki minus
this bag was $-0.010$ ($-0.023$ to 0.016) and $-0.012$ ($-0.031$ to
0.022) (donor bootstrap, 1,000 replicates). So whether MaxToki leads
depends on which gene order the bag uses. Most of global order
reflects lineage and cell-type grouping, which the structured null keeps
(null means 0.843 and 0.831, $p = 0.29$ and 0.15). S3~Appendix gives the
remaining checks, including a negative control that could not fail and
the single attention head singled out by the deployment, which beat
random projections only in trustworthiness.

\begin{figure}[H]
\caption{\textbf{MaxToki-217M: representation geometry after correction.}
(A)~Input gene-embedding tables: Pearson correlation between MaxToki's
gene--gene cosine-similarity matrix and that of Geneformer~V2-316M or
scGPT, on three gene panels from three tissue datasets (350--382 genes;
line, 95\% gene bootstrap interval). Dashed line, chance (genes
shuffled). (B)~Effective rank of the per-gene mean vectors at each layer,
for the trained model and three randomly initialised copies of the same
architecture, on the same 2,000 Tabula Sapiens immune cells and 1,500
genes (band, 95\% jackknife interval over 10 blocks of 200 cells).
(C)~Similarity (linear CKA) of the per-gene mean vectors between three
disjoint 2,000-cell samples, from layer~1, for the trained model and one
randomly initialised copy (band, 95\% jackknife interval over 30 groups
of 50 genes). (D)~Cross-lineage developmental order of cell-group
centroids on the external and zero-shot panels, for MaxToki-217M and for
three representations that use no model: the bag of each cell's input
tokens, a lookup of cell-type labels, and highly variable gene
expression (line, 95\% donor bootstrap interval, 2,000 replicates).}
\label{fig:caseGeom}
\end{figure}

\subsubsection*{Effective rank falls with depth in randomly initialised copies too}
We passed three disjoint samples of 2,000 Tabula Sapiens immune
cells~\cite{tabulasapiens2022} through MaxToki-217M and, at each layer,
averaged the hidden state of each of 1,500 genes over the cells that
contain it. Effective rank~\cite{roy2007effective} counts how many
directions these gene vectors use (at most 1,232, the hidden size). In the
trained model it was 561 at the token embedding (layer~0), 328 after the
first block and 227 at layer~11 (95\% jackknife interval over ten
200-cell blocks 224.8--229.5; 225.6 and 226.3 in the other two samples)
(Fig~\ref{fig:caseGeom}B). The deployed analysis, whose inputs were
encoded on the wrong scale, reported 94 at layer~11 for the same cells
and genes. Three randomly initialised copies of the architecture fell
from 816--817 at layer~0 to 390--421 at layer~11. From layer~1 to
layer~11 every random copy lost more of its effective rank than the
trained model (34.6--39.6\% against 30.8--31.3\%; the cell-block
jackknife intervals do not overlap), so the fall with depth does not need
training. Training lowers the effective rank at every layer, starting at
the embedding table, and adds two drops of its own: in the first block
(41.5\% against 21.0--21.3\% with random weights) and in the last step,
from layer~10 to layer~11, which runs the last block and then the final
normalisation (14.6--15.1\%, against a rise of 1.0--2.1\%; these two
parts were not measured separately). A feature-shuffle null cannot
separate these cases, because it stays far above the real effective rank
in every model.

Linear centred kernel alignment (CKA)~\cite{kornblith2019cka} measures how alike two sets of gene
vectors are, up to rotation and scale (1 means the same geometry).
Between the three samples it was 0.9997 at layer~1 and 0.9962 at
layer~11 in the trained model (95\% jackknife interval over 30 groups of
50 genes 0.9959--0.9964), against 0.9938 and 0.9647 (0.9598--0.9697) in a
randomly initialised copy (Fig~\ref{fig:caseGeom}C). In the trained model
most of the gap from 1 is sampling noise in the per-gene averages: with
1,000 instead of 2,000 cells per sample, $1-\mathrm{CKA}$ at layer~11
rose from 0.0038 to 0.0073. The trained model's gene averages are
therefore more stable across cell samples than a random model's, but
these values mainly measure how far a gene's average moves between random
sets of cells, and they should not be compared with CKA values reported
for other models. Other tissues, other initialisation schemes and
MaxToki-1B were not tested (S3~Appendix).

\subsubsection*{Cross-model alignment of gene embeddings}
This analysis compares the input gene-embedding tables of the three
models, which do not depend on how cells are encoded. We used three
overlapping panels of 350 to 382 genes, taken from three tissue
datasets: Tabula Sapiens lung, Tabula Sapiens immune cells, and an
external lung dataset. For each model we computed the cosine similarity
of every gene pair. We then took the Pearson correlation between two
models' similarity matrices. Between MaxToki-217M and scGPT this
correlation was 0.26--0.28. When genes were shuffled it was 0.00 (mean;
SD 0.004). We also report a held-out canonical correlation. Canonical
correlation measures how well linear projections of the two tables line
up. The held-out version is fitted on 80\% of the genes and scored on
the other 20\%, averaged over 10 random splits. Between MaxToki-217M and
scGPT it was 0.481--0.495, with chance about 0.00. MaxToki-217M agreed
more with Geneformer~V2-316M, which shares its gene vocabulary and token
identifiers: Pearson 0.39--0.40 and held-out canonical correlation
0.558--0.604 (Fig~\ref{fig:caseGeom}A).
The Geneformer-minus-scGPT difference in Pearson correlation was 0.10 to
0.14, with gene-bootstrap intervals that exclude zero on every panel. The
in-sample canonical correlation (fitted and scored on the same genes),
often used for such comparisons, is high even for unrelated tables. In
the deployed setting (each table reduced to 30 principal components, mean
of 10 canonical correlations) it was 0.41--0.43 for shuffled genes,
against 0.77--0.78 for Geneformer and 0.72--0.74 for scGPT. Its chance
level depends on the number of genes and of components. In the setting
the Study~A executors used (20 components, mean of all 20), chance was
about 0.20. So we report the held-out value. These are input tables only; they say nothing
about the models' internal representations.

\section*{Discussion}

\subsection*{What the evaluation supports}
The deployment produced many serious errors, and its own audit caught few
of them. The written sweep report of the in-session audit fully detected
10\% of the 68 reference errors and none of the 6 critical errors among
them. Across the full ledger, the report flagged none of the 11 critical
errors that were already present when the audit ran (one of them only
partly present). Later, during the audit's second repair round, the
deployment agent found one of these 11: the direction-of-change metric.
The other 2 of the 13 critical errors were made during the
audit's repair rounds. Fresh agents that audited the same snapshot found four times as
many reference errors, whether they used the checklist or a generic
instruction to review critically.

In the pre-registered comparisons with Claude Opus~5.5, the checklist did
not clearly help reviewers find more problems than a generic review of
the same budget in any of the three studies, and in none did it reduce
false alarms. There were two partial exceptions. In the exploratory arm
with Claude Haiku~4.5, a checklist review followed by repair raised the
share of pitfalls handled by 6.4 percentage points more than a generic
review followed by repair (95\% bootstrap interval over the 28 paired
executor runs, pooled over tasks, 1.3 to 11.7), but it did not raise the
share of correct verdicts. In Study~C the two checklist arms (generic and
as deployed) fully detected more of the 17 missing baselines and nulls
(51\% and 55\% against 39\% for the generic review) and of the 9 wrong
statistics (22\% and 37\% against 15\%), but fewer of the 12 code bugs
(69\% and 58\% against 81\%); these splits by error type were not tested
for significance. The specifications did not change whether a current
agent reached the correct conclusion on focused analysis tasks: with
Claude Opus~5.5 every executor run reached the correct verdict with or
without them. With a smaller model, Claude Haiku~4.5, they may help:
correct verdicts rose from 10 to 15 of 20 runs, a difference that 20 runs
per arm cannot confirm. The framework as deployed therefore did not make
the agent's work reliable, and the checklist added little over a generic
review of the same budget.

What did catch errors was independent review in a fresh context. The
blind re-audits of Study~C and the verification that built the ledger
were carried out by agents that had not produced the work and had no
stake in it. The in-session audit, by contrast, was written and applied
by the deployment agent itself, in one working session on the last day of
the deployment. Apart from a same-day extension of the topological
hypothesis screen, the work it audited had been done in other sessions.
The audit's sweep of all eight pipelines took about
8~minutes from request to report (3.7~hours counting the two repair
rounds), against 21 to 29~minutes for each blind re-audit, and it graded
run summaries more than code. We cannot separate the effects of
independence, more time and a different model: the deployment and its
in-session audit ran on Claude Opus~4.7, and the re-audits on Claude
Opus~5.5. Whatever the cause, an
audit by the agent that did the work was weak evidence that the work was
correct.

\subsection*{Most of the errors that mattered were in code}
Seven of the 13 critical errors were code bugs: an intervention hook that
edited the wrong tensor and deleted a transformer block (three ledger
entries), combined edits that overwrote each other, a model input encoded
on the wrong scale, a metric that ignored the prediction, and a lookup
that read the wrong embedding rows. None of these is a methodological
pattern of the kind the checklist describes, and the checklist did not
help fresh reviewers find them (the generic review found more of the code
bugs). Each of the seven, however, would have been caught by a short
deterministic test written for the analysis: an edit of zero must change
nothing; an edit at layer~$l$ must leave earlier layers unchanged;
combined edits must differ from single edits; the token order must match
the model's reference encoder on a few cells; a direction-of-change
accuracy must be compared with the constant-sign predictor on the same
records; and for a few named genes, the embedding row read must be the
row that the model's own vocabulary file gives them. Such tests are cheap
and do not depend on an agent's judgment. The other six critical errors
were not code bugs: three were missing or wrong nulls or baselines, two
were wrong statistics (one of them in the audit's own summary) and one
was an invalid proxy axis. The ledger places five of these six under a
checklist pattern (P1, P3 or P5 in Table~\ref{tab:checklist}; S1~Table).
The generic deterministic checks we built (structure of the
specification, parameter deviations, number tracing) catch a different
class of problem. By design they do not test whether a null randomises,
an edit lands on the right tensor, an input is encoded as the model
expects, a lookup reads the right rows or a number means what the text
says, so they would not have caught 12 of the 13 critical errors. The thirteenth, the audit's percentage with
no cell-by-cell record behind it, is the kind of problem the number
tracer looks for. We ran the tracer only on the run summaries, not on the
audit reports, so we did not test whether it would have flagged it. For
agent-run interpretability we would therefore recommend
analysis-specific tests of interventions, inputs, lookups and baselines,
run before results are summarised, together with independent
fresh-context review, rather than a longer checklist.

The held-out benchmark adds a second observation. Every review of a
flawed, single-analysis package found its documented error. This
included a raw-count input fed to an encoder that expects binned values,
the same kind of error as MaxToki's input encoding. Only part of
the MaxToki input-encoding error was in the Study~C reference set: three attention runs
fed $\log(1+\text{CP10k})$ values in as if they were counts (S1~Table,
L118). Of the nine blind audits of the whole MaxToki project, three were
graded as fully detecting this part and two as detecting it partly. These
five suspected that the Adamson input was log-normalised and noted that
no code checked for raw counts. None found that every run except the RPE1
attention run was encoded on the wrong scale.
Reviewing one analysis at a time, with its code, outputs and summary
together, appears far more effective than reviewing a large project in
one pass. This comparison is across studies, with different budgets per
analysis, and should be tested directly.

Reviews also over-called errors. Most reviews of correct analyses (53 of
66) claimed at least one error in a main result that was not one. On
correct and flawed analyses together, the reviews made 105 such claims.
The graders marked 82 of them as false alarms (73 on correct and 9 on
flawed analyses); the verifier agent ruled the other 23 false alarms
after a grader dispute. A post-hoc check of the 82, which was not
pre-registered, found that 76 were accurate observations whose severity
was overstated. A review step
is therefore only as good as the step that checks its points against the
data before acting on them. In Study~A a repair step was meant to do
this, and it did not always succeed. With Opus~5.5, repairs after generic
reviews turned 3 of 40 correct verdicts into wrong ones, and repairs
after checklist reviews turned none. In all three cases the reviewer
leaned on metrics named in the specification, one of which has a high
chance level that the specification does not state. A specification can
thus mislead a reviewer even when the executor that followed it was not
misled. In the exploratory Haiku~4.5 arm the pattern was reversed:
repairs after checklist reviews turned 3 of 21 correct verdicts wrong
(and 3 wrong verdicts right), and repairs after generic reviews turned
none wrong. So each review style, with one of the two models, led
repairs to turn correct verdicts wrong.

\subsection*{Specifications and prompts}
During the deployment a specification was, in practice, a long and
structured prompt: nothing checked that it was complete, that its
parameters were used or that its validation criteria were met, and the
executing agent changed parameters (for example, the number of null
draws) without recording why in most cases. What can be checked by code
is narrow (Table~\ref{tab:checks}): document structure, resolvable code
references, deviations from machine-readable parameters, and whether a
stored gate flag matches the stored numbers. Whether a null randomises,
an intervention edits the intended tensor or a number means what the text
says remains a matter of judgment unless an analysis-specific test is
written for it. The specifications also carry expectations from earlier
studies; one stated that a negative verdict was the expected outcome. For
a strong agent on a focused task this made no difference to the verdict,
but a contract that states the expected answer is a source of bias that a
pre-registered analysis plan would avoid.

\subsection*{Transfer to other models and settings}
The held-out items came from other projects by the author: a text-model
cell embedding (Cell2Sentence), a protein language model (ESM-2), published
single-cell transformers (scGPT, Geneformer and STATE), and
gene-regulatory-network benchmarks that used no model. Every documented
error was found in each of these four groups. In three groups most
reviews of correct analyses raised a false alarm: 18 of 18 reviews for
Cell2Sentence, 12 of 18 for ESM-2 and 17 of 18 for single-cell
transformers. In the gene-regulatory-network group half did (6 of 12).
Each group had only two or three correct items, too few to compare the
groups. Because the checklist's patterns concern study
design rather than any model's internals, we expect the review findings
to transfer to other foundation models in biology. The code-level tests do not transfer
unchanged: each architecture and tokenizer needs its own tests of where an
edit lands and how inputs are encoded. We tested one agent model family
(Claude) as executor, reviewer and grader, one deployment, and analyses
from one author's projects; other agent families, other research groups'
analyses and human reviewers were not tested.

\subsection*{No human or non-agent baseline}
We did not compare the agent workflow with human analysts or human
reviewers, and we make no claim about efficiency. The deployment's
supervision time was not logged, and the controlled studies measured
agent effort (tool calls, time), not human effort. Whether a human
reviewer with the same budget would find more of the code-level errors,
or raise fewer false alarms, is an open question that a study with human
reviewers could answer on the same items, which are released for this
purpose.

\subsection*{What the case study says about MaxToki-217M, and what it does not}
The corrected analyses describe MaxToki-217M given one cell at a time,
with rank-value inputs from K562, RPE1 and Tabula Sapiens cells, and
judged against TRRUST, DoRothEA and CRISPR-interference screens. They do
not test MaxToki's multi-cell trajectory mode, the larger MaxToki-1B
(except for one attention run, whose inputs were encoded incorrectly and
which we do not report), or other single-cell foundation models. Negative
results here mean that these probes, with these references, found no
signal beyond simple baselines in this model; they are not evidence that
\scfm{}s in general lack regulatory information. The references themselves
are limited: TRRUST covers few transcription factors per cell line, and
the knockdown fold-changes are noisy, so modest signals could be missed.

\subsection*{Limitations}
The executor, reviewer, repairer and grader agents all came from one model
family, so shared blind spots cannot be excluded; graders agreed closely
with each other, but no human graded the deliverables. Study~A showed a
ceiling for Opus~5.5, which leaves it without power to detect benefits of
the specification for that model; the exploratory Haiku arm was added
after that ceiling was seen and has 20 runs per arm (S2~Appendix). The Study~A tasks were built from the
deployment's data and so are not independent of it. The Study~B items
came from the author's own projects. Each flawed item held one documented
error, and every review found it, in both arms (33 of 33 per arm). So
Study~B could not show whether the checklist finds more of the documented
errors. It could compare only false alarms and the number of other real
problems found per review (2.1 in each arm). The error ledger is a lower bound:
blind re-audits added 60 errors to it and the full input-encoding error was
found last, so detection rates are relative to the errors known when
each study's key was frozen. The deployment ran on Claude Opus~4.7 (its
one automated loop called Claude Sonnet), while the evaluation used
Opus~5.5 and, in the exploratory arm, Haiku~4.5. So a difference in
model may explain part of the gap between the in-session audit and the
blind re-audits. Finally, this manuscript
was drafted by an agent under the author's direction. The source data of
every figure, the error ledger with the file location of each error
(S1~Table), and a results map that lists every number in the paper with
its source file (RESULTS\_MAP.csv) are released with the code.

\section*{Conclusion}
An agent following written method specifications ran eight
interpretability analyses of MaxToki-217M. The report of its own
checklist audit, written before any repair, caught 10\% of a fixed set of
68 errors in these analyses. It caught none of the 11 critical errors
present when the audit ran. Fresh agents auditing the same work caught
41--45\% of the 68, with or without the checklist. In the pre-registered
tests with Claude Opus~5.5, every executor run reached the correct
conclusion with or without the specifications (20 of 20 in each arm), so
these tests could not show a benefit. The checklist did not clearly find
more errors than a generic review of the same budget, and it did not
raise fewer false alarms. There were two exceptions in our data. In an
exploratory arm with the smaller Claude Haiku~4.5, a checklist review
followed by repair raised the share of pitfalls handled by 6.4
percentage points more than a generic review followed by repair (95\%
bootstrap interval over 28 paired executor runs pooled over tasks, 1.3
to 11.7), but it did not raise the number of correct verdicts. In the
blind re-audits, the two checklist arms (generic and as deployed) fully
detected more of the 17 missing baselines and nulls than the generic
review (51\% and 55\% against 39\%) and more of the 9 wrong statistics
(22\% and 37\% against 15\%). They fully detected fewer of the 12 code
bugs (69\% and 58\% against 81\%). These splits by error type were not
tested for significance. For agent-run interpretability we therefore
recommend independent review in a fresh context, short code tests of
interventions, inputs and baselines, and a check of every review point
against the data before it is acted on. After correction, the attention
scores, sparse-autoencoder features and traced circuits of MaxToki-217M
did not beat gene-level baselines with the references used here. One
pair-level signal remained. In RPE1 cells, symmetric attention beat
rewired TRRUST networks that keep every gene's number of links, even
after the gene-level features of both genes were removed ($p = 0.003$,
not adjusted for multiple tests). Co-expression from the same cells
carried a signal of about the same size. This signal did not replicate
in K562 cells.

\section*{Supporting information}

\paragraph*{S1 Table.}
\label{S1_Table}
{\bf Error ledger.} The 149 confirmed errors in the deployment's code,
logs, outputs, run summaries and audit reports, with location,
confirming evidence, one primary type (with a note when more than one
type applies), checklist pattern, severity, presence before the
in-session audit, and who or what first found each one (CSV). For the errors
found in the blind re-audits, the confirming code was not kept; their
rows give the recomputed values.

\paragraph*{S2 Appendix.}
\label{S2_Appendix}
{\bf Protocol of the controlled studies.} The pre-registered design,
outcomes and analyses; the three amendments (the first logged after
Study~C and before Studies~A and~B, the other two for the exploratory
arm); the freeze log with file hashes; the generic and deployed
checklists; the prompts sent to the agents; the Study~A task briefs,
answer keys and grading rubrics; the Study~B items and detection rules;
the Study~C reference set; and the exploratory arm with Claude Haiku~4.5
as executor, reviewer and repairer.

\paragraph*{S3 Appendix.}
\label{S3_Appendix}
{\bf Corrected MaxToki-217M analyses.} Methods, checks and full results
of each corrected analysis. The key numbers of each corrected analysis
were computed a second way with separate code. Other agents also
re-derived them with their own code: three for the RPE1 attention
analysis and one for the cross-model comparison.

\section*{Acknowledgments}
The author thanks the developers and maintainers of the public data
resources and model releases on which this study depends.

\nolinenumbers

\bibliography{references}

\end{document}
